% Proofs for the section on approximate meta-agents and controlled renormalization.
% Every proof below was reconstructed from the hypotheses of the statements by the independent check.
\section{Proofs for \texorpdfstring{Section~\ref{sec:apx}}{approximate meta-agents}}\label{app:apx}

Throughout, $\lvert v\rvert$ is the Euclidean norm, $\Vert M\Vert$ the spectral norm, $\Vert M\Vert_{\mathrm F}$ the Frobenius norm, $\Vert X\Vert_{L^2}=(\E\lvert X\rvert^2)^{1/2}$, and $\preceq$ the order of symmetric matrices. The matrices $C^\sharp$ and $Z^\sharp$ are the blocks of the inverse of $\binom CZ$, so that $CC^\sharp=I$, $ZZ^\sharp=I$, $CZ^\sharp=0$, $ZC^\sharp=0$ and $C^\sharp C+Z^\sharp Z=I$, and $A_{yy}=CAC^\sharp$, $A_{y\perp}=CAZ^\sharp$, $A_{\perp y}=ZAC^\sharp$, $A_{\perp\perp}=ZAZ^\sharp$. The randomization $\omega$ of a causal coarse policy takes values in a standard Borel space, has the same law under the fine process and under the declared model, and is independent of $(x_0,\mathfrak n)$ in the first and of $(\hat y_0,\hat w)$ in the second; the second moments $\E\lvert u_t\rvert^2$ are finite under the fine process. Regular conditional laws exist because all spaces are standard Borel. For a symmetric $S\succeq0$ with $CSC^\top\succ0$ put
\begin{equation}\label{eq:apx-condition-map}
\Pi_C(S)=S-SC^\top(CSC^\top)^{-1}CS ,
\end{equation}
the covariance of a Gaussian vector with covariance $S$ given its image under $C$; it satisfies $C\Pi_C(S)=0$ and $0\preceq\Pi_C(S)\preceq S$, and it is nondecreasing in $S$, because $\Pi_C(S)=\min_K(I-KC)S(I-KC)^\top$.

\begin{lemma}[Separation]\label{lem:apx-separation}
\begin{enumerate}
\item[(a)] For every causal coarse policy and every $t\ge0$, the conditional law of $x_t$ given $\mathcal G_t=\sigma(y_{0:t},\omega)$ is $\Ncal(\hat x_t,P_t)$, where $\hat x_t$ is a measurable function of $(y_{0:t},\omega)$ with $C\hat x_t=y_t$ and
\begin{equation}\label{eq:apx-pt}
P_0=\Pi_C(\Sigma_0),\qquad P_{t+1}=\Pi_C\bigl(AP_tA^\top+Q\bigr)
\end{equation}
is deterministic and the same for all causal coarse policies, with $CP_t=0$. The conditional law of $y_{t+1}$ given $\mathcal G_t$ is $\Ncal(CA\hat x_t+CBu_t,S_t)$ with $S_t=Q^y+CAP_tA^\top C^\top\succ0$.
\item[(b)] Let $s_{t+1}=Hy_t+v_t$ with $v_t\sim\Ncal(0,R)$ independent of each other and of $(x_0,\mathfrak n,\omega)$, $R\succ0$, and let $u_t=\pi_t(s_{1:t},\omega)$. Then the conditional laws of $x_t$ given $(s_{1:t},\omega)$ and given $(s_{1:t+1},\omega)$ are $\Ncal(\hat x^-_t,\Sigma^-_t)$ and $\Ncal(\hat x^{\mathrm f}_t,\Sigma^{\mathrm f}_t)$ with
\begin{equation}\label{eq:apx-record-filter}
\Sigma^-_0=\Sigma_0,\qquad\Sigma^{\mathrm f}_t=\Sigma^-_t-\Sigma^-_tC^\top H^\top\bigl(HC\Sigma^-_tC^\top H^\top+R\bigr)^{-1}HC\Sigma^-_t,\qquad\Sigma^-_{t+1}=A\Sigma^{\mathrm f}_tA^\top+Q ,
\end{equation}
which no such policy changes, and $\hat x^-_{t+1}=A\hat x^{\mathrm f}_t+Bu_t$.
\end{enumerate}
\end{lemma}

\begin{proof}
(a) The proof is by induction on $t$. Since $x_0\sim\Ncal(m_0,\Sigma_0)$ is independent of $\omega$ and $C\Sigma_0C^\top\succ0$, Gaussian conditioning on $y_0=Cx_0$ gives the law $\Ncal(m_0+\Sigma_0C^\top(C\Sigma_0C^\top)^{-1}(y_0-Cm_0),\Pi_C(\Sigma_0))$, and conditioning on $\omega$ as well changes nothing. Assume the claim at $t$. By construction $x_t$ and $y_{0:t}$ are measurable functions of $(x_0,\mathfrak n_{0:t-1},\omega)$, so $\mathfrak n_t$ is independent of $(x_t,\mathcal G_t)$, and $u_t$ is $\mathcal G_t$-measurable. Given $\mathcal G_t$ the vector $x_{t+1}=Ax_t+Bu_t+\mathfrak n_t$ is therefore Gaussian with mean $A\hat x_t+Bu_t$ and covariance $AP_tA^\top+Q$, and its image $y_{t+1}=Cx_{t+1}$ is Gaussian with mean $CA\hat x_t+CBu_t$ and covariance $C(AP_tA^\top+Q)C^\top=S_t\succeq Q^y\succ0$. Conditioning the pair $(x_{t+1},y_{t+1})$, which is jointly Gaussian given $\mathcal G_t$, on $y_{t+1}$ gives a Gaussian law with covariance $\Pi_C(AP_tA^\top+Q)$ and a mean that is an affine function of $y_{t+1}$ with $\mathcal G_t$-measurable coefficients. Since $\mathcal G_{t+1}=\mathcal G_t\vee\sigma(y_{t+1})$, this is the claim at $t+1$. The identities $C\hat x_t=y_t$ and $CP_t=0$ hold because $Cx_t=y_t$ is $\mathcal G_t$-measurable.

(b) The same induction applies with the two conditioning steps of \cref{eq:apx-record-filter}. Given $(s_{1:t},\omega)$ and the induction hypothesis $x_t\sim\Ncal(\hat x^-_t,\Sigma^-_t)$, the reading $s_{t+1}=HCx_t+v_t$ is jointly Gaussian with $x_t$, because $v_t$ is independent of $(x_t,s_{1:t},\omega)$, and conditioning on it gives $\Ncal(\hat x^{\mathrm f}_t,\Sigma^{\mathrm f}_t)$. The act $u_t$ is measurable in $(s_{1:t},\omega)$ and $\mathfrak n_t$ is independent of $(x_t,s_{1:t+1},\omega)$, so $x_{t+1}$ given $(s_{1:t+1},\omega)$ is $\Ncal(A\hat x^{\mathrm f}_t+Bu_t,A\Sigma^{\mathrm f}_tA^\top+Q)$.
\end{proof}

\begin{remark}[Starts that depend on the randomization]\label{rem:apx-separation-start}
Both parts of \cref{lem:apx-separation} hold, with the same proof, when $x_0$ is not independent of $\omega$ but its law given $\omega$ is $\Ncal(m_0,\Sigma_0)$ with a deterministic $\Sigma_0\succeq0$ and a measurable function $m_0$ of $\omega$ with $\E\lvert m_0\rvert^2<\infty$, provided the noise, $\mathfrak n$ and in (b) also $v$, is independent of $(x_0,\omega)$. Part (a) needs $C\Sigma_0C^\top\succ0$ and part (b) does not. The induction starts from this conditional law, with $\hat x^-_0=m_0$ in (b), and each later step uses only that $v_t$ is independent of $(x_0,\mathfrak n_{0:t-1},v_{0:t-1},\omega)$ and $\mathfrak n_t$ of $(x_0,\mathfrak n_{0:t-1},v_{0:t},\omega)$. The recursions \cref{eq:apx-pt} and \cref{eq:apx-record-filter} are unchanged, with $\Sigma_0$ the conditional covariance.
\end{remark}

\begin{proof}[Proof of \cref{thm:apx-path}]
Both $P^\pi$ and $\hat P^\pi$ are laws of $(\omega,y_0,\dots,y_T)$ on the same standard Borel space. Under $P^\pi$ the randomization has its law, $y_0\sim P_{y_0}=\Ncal(Cm_0,C\Sigma_0C^\top)$ is independent of it, and by \cref{lem:apx-separation}(a) the kernel of $y_{t+1}$ given $(\omega,y_{0:t})$ is $\Ncal(CA\hat x_t+CBu_t,S_t)$. Under $\hat P^\pi$ the randomization has the same law, $y_0\sim\hat P_0$ is independent of it, and the kernel of $y_{t+1}$ given $(\omega,y_{0:t})$ is $\Ncal(\bar Ay_t+\bar Bu_t,\bar Q)$ with the same $u_t=\pi_t(y_{0:t},\omega)$. The chain rule for relative entropy, which holds with values in $[0,\infty]$ on standard Borel spaces, gives
\begin{equation*}
\KL\bigl(P^\pi\Vert\hat P^\pi\bigr)=\KL\bigl(P_{y_0}\Vert\hat P_0\bigr)+\sum_{t=0}^{T-1}\E^\pi\KL\Bigl(\Ncal\bigl(CA\hat x_t+CBu_t,S_t\bigr)\Big\Vert\Ncal\bigl(\bar Ay_t+\bar Bu_t,\bar Q\bigr)\Bigr),
\end{equation*}
the term of $\omega$ being zero. The relative entropy of $\Ncal(m,S)$ from $\Ncal(\hat m,\bar Q)$ is $h(\bar Q^{-1}S)+\frac12\Vert m-\hat m\Vert^2_{\bar Q^{-1}}$, and $CA\hat x_t-\bar Ay_t=(CA-\bar AC)\hat x_t=D\hat x_t$ because $y_t=C\hat x_t$. Hence the summand is $h(\bar Q^{-1}S_t)+\frac12\E^\pi\Vert D\hat x_t+Eu_t\Vert^2_{\bar Q^{-1}}$, the second form of the statement; it is finite because $\E\lvert u_t\rvert^2<\infty$ and, by induction on the recursion, $\E\lvert x_t\rvert^2<\infty$.

For the first form write $\delta_t=D\hat x_t+Eu_t+D(x_t-\hat x_t)$. Given $\mathcal G_t$ the last term has mean zero and covariance $DP_tD^\top$, and the first two are $\mathcal G_t$-measurable, so $\E^\pi\Vert\delta_t\Vert^2_{\bar Q^{-1}}=\E^\pi\Vert D\hat x_t+Eu_t\Vert^2_{\bar Q^{-1}}+\tr(\bar Q^{-1}DP_tD^\top)$. Since $CP_t=0$, $DP_tD^\top=CAP_tA^\top C^\top=S_t-Q^y$. Finally
\begin{equation*}
h\bigl(\bar Q^{-1}S_t\bigr)=h\bigl(\bar Q^{-1}Q^y\bigr)+\tfrac12\tr\bigl(\bar Q^{-1}CAP_tA^\top C^\top\bigr)-\tfrac12\ln\det\bigl((Q^y)^{-1}S_t\bigr),
\end{equation*}
and the term subtracted on the right is $\mathfrak d_{\mathrm{dyn},t}$. Adding the two displays gives \cref{eq:apx-path}. The expression for $\mathfrak d_{\mathrm{dyn},t}$ is that of \cref{prop:met-closure-notions}(b) with $\mathsf P=C$, $Q^{\mathfrak n}=Q$ and $P^{\mathrm f}_t=P_t$; by \cref{lem:apx-separation}(a) it is also the conditional mutual information $I(y_{t+1};x_t\mid y_{0:t},\omega)$ under every causal coarse policy.
\end{proof}

\begin{proof}[Proof of \cref{cor:apx-uniform}]
The moment bounds of the statement are
\begin{equation}\label{eq:apx-moments}
\begin{aligned}
\bar X&=\frac{1}{1-g_\kappa}\Bigl[M\mu_0+\frac{M\Vert B\Vert\bar c}{1-\alpha}+M\Bigl(\frac{\tr Q}{1-\alpha^2}\Bigr)^{1/2}\Bigr],\\
\bar X_\perp&=\min\Bigl\{\Vert Z\Vert\bar X,\ M_\perp\Vert Z\Vert\mu_0+\frac{M_\perp}{1-\alpha_\perp}\bigl[(\Vert A_{\perp y}\Vert+\kappa\Vert B_\perp\Vert)\Vert C\Vert\bar X+\Vert B_\perp\Vert\bar c\bigr]+M_\perp\Bigl(\frac{\tr Q^\perp}{1-\alpha_\perp^2}\Bigr)^{1/2}\Bigr\},
\end{aligned}
\end{equation}
and the claim is $\sup_t\Vert x_t\Vert_{L^2}\le\bar X$ and $\sup_t\Vert x_{\perp,t}\Vert_{L^2}\le\bar X_\perp$ for every $\pi\in\mathcal P(\kappa,\bar c)$. The first entry of the minimum does not use the hypothesis on $A_{\perp\perp}$, which can therefore be dropped with $\bar X_\perp=\Vert Z\Vert\bar X$.

Let $a_t=\Vert x_t\Vert_{L^2}$. From $x_t=A^tx_0+\sum_{s<t}A^{t-1-s}(Bu_s+\mathfrak n_s)$, the bound $\lvert u_s\rvert\le\kappa\Vert C\Vert\lvert x_s\rvert+\bar c$, the independence of the noise terms and $\tr(A^jQA^{j\top})\le\Vert A^j\Vert^2\tr Q$,
\begin{equation*}
a_t\le M\mu_0+\frac{M\Vert B\Vert}{1-\alpha}\Bigl(\kappa\Vert C\Vert\max_{s<t}a_s+\bar c\Bigr)+M\Bigl(\frac{\tr Q}{1-\alpha^2}\Bigr)^{1/2}=(1-g_\kappa)\bar X+g_\kappa\max_{s<t}a_s .
\end{equation*}
Since $M\ge1$ and $g_\kappa<1$, $a_0\le\mu_0\le\bar X$, and induction gives $a_t\le\bar X$ for all $t$. Hence $\Vert x_{\perp,t}\Vert_{L^2}\le\Vert Z\Vert\bar X$. The second entry of the minimum follows in the same way from $x_{\perp,t+1}=A_{\perp\perp}x_{\perp,t}+A_{\perp y}y_t+B_\perp u_t+Z\mathfrak n_t$, with $\Vert y_s\Vert_{L^2}\le\Vert C\Vert\bar X$ and $\Vert u_s\Vert_{L^2}\le\kappa\Vert C\Vert\bar X+\bar c$.

Since $x=C^\sharp y+Z^\sharp x_\perp$, $Dx=(A_{yy}-\bar A)y+A_{y\perp}x_\perp$, so $\Vert\delta_t\Vert_{L^2}\le\Vert A_{yy}-\bar A\Vert\Vert C\Vert\bar X+\Vert A_{y\perp}\Vert\bar X_\perp+\Vert E\Vert(\kappa\Vert C\Vert\bar X+\bar c)=\bar{\mathfrak d}$. Inserting $\E^\pi\Vert\delta_t\Vert^2_{\bar Q^{-1}}\le\bar{\mathfrak d}^2/\lambda_{\min}(\bar Q)$ and $\mathfrak d_{\mathrm{dyn},t}\ge0$ in \cref{eq:apx-path} gives \cref{eq:apx-uniform}.
\end{proof}

\begin{remark}[The small-gain condition]\label{rem:apx-smallgain}
Let $n=2$, $A=\mathsf aI$ with $0<\mathsf a<1$, $Q=I$, $C=(1\ \ 0)$, $B=(1\ \ 0)^\top$, $\bar A=\mathsf a$, $\bar B=\frac12$ and $u_t=\kappa y_t$ with $\kappa>0$. Then $D=0$, $E=\frac12$, $\delta_t=\frac12\kappa y_t$, $M=1$, $\alpha=\mathsf a$ and $g_\kappa=\kappa/(1-\mathsf a)$. The closed loop is $y_{t+1}=(\mathsf a+\kappa)y_t+\mathfrak n_{1,t}$, so $\E\delta_t^2=\frac14\kappa^2[(\mathsf a+\kappa)^{2t}\E y_0^2+\sum_{j<t}(\mathsf a+\kappa)^{2j}]$, which is bounded in $t$ exactly when $g_\kappa<1$ and grows like $(\mathsf a+\kappa)^{2t}$ when $g_\kappa>1$. In this family the condition $g_\kappa<1$ of \cref{cor:apx-uniform} is therefore necessary for a bound linear in $T$. In general it is sufficient and not necessary, since it ignores the sign of the feedback.
\end{remark}

\begin{proof}[Proof of \cref{prop:apx-timestep}]
With the stated scaling, $D=CA-\bar AC=\mathsf hD_{\mathrm c}$, $E=\mathsf hE_{\mathrm c}$, $Q^y=\mathsf h\mathsf q^y$ and $\bar Q^{-1}Q^y=\bar{\mathsf q}^{-1}\mathsf q^y$, and the horizons are the integers $T=\tau/\mathsf h$. Because $CP_t=0$, $CAP_tA^\top C^\top=\mathsf h^2CFP_tF^\top C^\top$, so that $\bar Q^{-1}S_t=\bar{\mathsf q}^{-1}(\mathsf q^y+\mathsf hCFP_tF^\top C^\top)$.

Lower bound. By the second form of the summand in \cref{thm:apx-path}, $\KL(P^\pi\Vert\hat P^\pi)\ge\sum_{t<T}h(\bar Q^{-1}S_t)$ for every causal coarse policy. By \cref{eq:apx-pt} and the monotonicity of $\Pi_C$, $P_t$ is at most the covariance of the uncontrolled state, so for $t\le T$
\begin{equation*}
\Vert P_t\Vert\le\Vert A\Vert^{2t}\Vert\Sigma_0\Vert+\mathsf h\Vert\mathsf q\Vert\sum_{j<t}\Vert A\Vert^{2j}\le e^{2\Vert F\Vert\tau}\bigl(\Vert\Sigma_0\Vert+\tau\Vert\mathsf q\Vert\bigr),
\end{equation*}
because $\Vert A\Vert\le e^{\mathsf h\Vert F\Vert}$. Hence $\bar Q^{-1}S_t\to\bar{\mathsf q}^{-1}\mathsf q^y$ uniformly in $t\le T$ as $\mathsf h\to0$, and since $h$ is continuous on matrices with positive determinant, $\mathsf h\sum_{t<T}h(\bar Q^{-1}S_t)\to\tau h(\bar{\mathsf q}^{-1}\mathsf q^y)$. The limit is positive when $\bar{\mathsf q}\ne\mathsf q^y$, because $h(M)=\frac12\sum_ig(\lambda_i)$ over the eigenvalues of $M$, which are positive here, and $g$ vanishes only at $1$. The bound holds uniformly over the causal coarse policies, which may depend on $\mathsf h$.

Upper bound. When $\bar{\mathsf q}=\mathsf q^y$ the first term of the summand in \cref{eq:apx-path} vanishes, and $\mathfrak d_{\mathrm{dyn},t}\ge0$, so the summand is at most $\frac12\E^\pi\Vert\mathsf hD_{\mathrm c}x_t+\mathsf hE_{\mathrm c}u_t\Vert^2_{(\mathsf h\mathsf q^y)^{-1}}=\frac{\mathsf h}{2}\E^\pi\Vert D_{\mathrm c}x_t+E_{\mathrm c}u_t\Vert^2_{(\mathsf q^y)^{-1}}$. For the boundedness let $a_t=\Vert x_t\Vert_{L^2}$, $\mathsf c_1=\Vert F\Vert+\kappa\Vert G\Vert\Vert C\Vert$ and $\mathsf h\le1$. Since $\mathfrak n_t$ is independent of $(x_t,u_t)$ and $\lvert u_t\rvert\le\kappa\Vert C\Vert\lvert x_t\rvert+\bar c$,
\begin{equation*}
a_{t+1}^2\le\bigl((1+\mathsf h\mathsf c_1)a_t+\mathsf h\Vert G\Vert\bar c\bigr)^2+\mathsf h\tr\mathsf q\le(1+\mathsf h)(1+\mathsf h\mathsf c_1)^2a_t^2+\mathsf h\bigl[(1+\mathsf h)\Vert G\Vert^2\bar c^2+\tr\mathsf q\bigr],
\end{equation*}
by $(p+q)^2\le(1+\mathsf h)p^2+(1+\mathsf h^{-1})q^2$. With $(1+\mathsf h)(1+\mathsf h\mathsf c_1)^2\le e^{\mathsf h(1+2\mathsf c_1)}$ this gives, for all $t\le T$,
\begin{equation*}
a_t^2\le\bar X_\tau^2:=e^{(1+2\mathsf c_1)\tau}\bigl(\mu_0^2+\tau\bigl[2\Vert G\Vert^2\bar c^2+\tr\mathsf q\bigr]\bigr),
\end{equation*}
and therefore
\begin{equation*}
\KL\bigl(P^\pi\Vert\hat P^\pi\bigr)-\KL\bigl(P_{y_0}\Vert\hat P_0\bigr)\le\frac{\tau}{2\lambda_{\min}(\mathsf q^y)}\Bigl(\Vert D_{\mathrm c}\Vert\bar X_\tau+\Vert E_{\mathrm c}\Vert\bigl(\kappa\Vert C\Vert\bar X_\tau+\bar c\bigr)\Bigr)^2
\end{equation*}
for every $\pi\in\mathcal P(\kappa,\bar c)$ and every $\mathsf h\le1$, a bound that does not depend on $\mathsf h$. No stability or small-gain hypothesis is used on a fixed physical horizon.
\end{proof}

\begin{proof}[Proof of \cref{prop:apx-kernel-floor}]
Here $C=(1\ \ 0)$, $Z=(0\ \ 1)$, $A_{yy}=\mathsf b$, $A_{y\perp}=\mathsf g$, $A_{\perp y}=0$ and $A_{\perp\perp}=\mathsf a$, so $\Gamma_l=A_{y\perp}A_{\perp\perp}^{l-1}A_{\perp y}=0$ for every $l\ge1$, and the orthogonal force of \cref{thm:agy-mori-zwanzig} is $\mathfrak f_{t+1}=\mathsf gx_{2,t}+\mathfrak n_{1,t}$, of stationary variance $1+\mathsf g^2V$ and of covariance $\mathsf g^2\mathsf a^jV$ at lag $j\ge1$. Truncating the kernel at any order and replacing $\mathfrak f$ by white Gaussian noise of variance $\sigma^2>0$ gives the declared model $\bar A=\mathsf b$, $\bar Q=\sigma^2$, without input, for which $D=CA-\bar AC=(0\ \ \mathsf g)$.

Let $x_0\sim\Ncal(0,\Sigma_0)$ with $\Var(x_{2,0})=V$, so that $x_{2,t}$ is stationary with mean zero. By \cref{lem:apx-separation}(a), $P_t$ has the single nonzero entry $P^\perp_t=\Var(x_{2,t}\mid y_{0:t})$, and \cref{eq:apx-pt} reads
\begin{equation*}
P^\perp_{t+1}=f\bigl(P^\perp_t\bigr),\qquad f(P)=1+\frac{\mathsf a^2P}{1+\mathsf g^2P},
\end{equation*}
while $S_t=1+\mathsf g^2P^\perp_t$. The map $f$ is nondecreasing, concave and bounded on $[0,\infty)$ with $f(0)=1$, so $f(P)-P$ has a single zero $P_\infty>0$, is positive below it and negative above it, and $P^\perp_t$ converges monotonically to $P_\infty$. Since $V=\mathsf a^2V+1$, $f(V)-V=-\mathsf a^2\mathsf g^2V^2/(1+\mathsf g^2V)$, which is negative when $\mathsf a\ne0$ and zero when $\mathsf a=0$; hence $P_\infty<V$ if $\mathsf a\ne0$ and $P_\infty=V=1$ if $\mathsf a=0$.

By the second form of the summand in \cref{thm:apx-path} and $\E\hat x_{2,t}^2=V-P^\perp_t$, the summand at time $t$ is
\begin{equation*}
h\bigl(S_t/\sigma^2\bigr)+\frac{\mathsf g^2(V-P^\perp_t)}{2\sigma^2}=\frac12\Bigl[\frac{1+\mathsf g^2V}{\sigma^2}-1-\ln\frac{1+\mathsf g^2P^\perp_t}{\sigma^2}\Bigr].
\end{equation*}
The initial term of \cref{eq:apx-path} is finite, so the rate $\lim_{T\to\infty}T^{-1}\KL(P\Vert\hat P)$ exists and equals the right side with $P_\infty$ in place of $P^\perp_t$. As a function of $\sigma^2$ this is minimal at $\sigma^2=1+\mathsf g^2V$, where it equals $\frac12\ln((1+\mathsf g^2V)/(1+\mathsf g^2P_\infty))$. This is positive when $\mathsf a\ne0$. When $\mathsf a=0$ the force $\mathfrak f$ is white, the coarse process is a first-order autoregression and the rate at $\sigma^2=1+\mathsf g^2$ is zero.
\end{proof}

\begin{remark}[Scope of \cref{prop:apx-kernel-floor}]\label{rem:apx-kernel-scope}
The kernels of \cref{thm:agy-mori-zwanzig} depend on the complement $Z$. The proposition concerns the complement $x_\perp=x_2$, for which $A_{\perp y}=0$. For the complement of part (a) of that theorem, chosen so that $x_\perp$ and $y$ are uncorrelated in the stationary law, the kernels of this example do not vanish in general, and the truncation of those kernels is not analyzed here.
\end{remark}

The filter of the unresolved coordinate, used in \cref{thm:apx-order-d}, is the following. Let $\hat\xi_t=Z\hat x_t=\E[x_{\perp,t}\mid\mathcal G_t]$ and $P^\perp_t=ZP_tZ^\top$. Since $CP_t=0$, $P_t=Z^\sharp P^\perp_tZ^{\sharp\top}$, so $S_t=Q^y+A_{y\perp}P^\perp_tA_{y\perp}^\top$ and $CA\hat x_t=A_{yy}y_t+A_{y\perp}\hat\xi_t$. Given $\mathcal G_t$, the pair $(x_{\perp,t+1},y_{t+1})$ is Gaussian with cross-covariance $A_{\perp\perp}P^\perp_tA_{y\perp}^\top+Q^{\perp y}$, and conditioning on $y_{t+1}$ gives, with
\begin{equation*}
G_t=\bigl(A_{\perp\perp}P^\perp_tA_{y\perp}^\top+Q^{\perp y}\bigr)S_t^{-1},\qquad\Phi_t=A_{\perp\perp}-G_tA_{y\perp},
\end{equation*}
the recursion
\begin{equation}\label{eq:apx-perp-filter}
\hat\xi_{t+1}=\Phi_t\hat\xi_t+A_{\perp y}y_t+B_\perp u_t+G_t\bigl(y_{t+1}-A_{yy}y_t-B_yu_t\bigr).
\end{equation}
The error $\epsilon_t=x_{\perp,t}-\hat\xi_t$ obeys $\epsilon_{t+1}=\Phi_t\epsilon_t+Z\mathfrak n_t-G_tC\mathfrak n_t$, with $\mathfrak n_t$ independent of $\epsilon_t$, so
\begin{equation}\label{eq:apx-joseph}
P^\perp_{t+1}=\Phi_tP^\perp_t\Phi_t^\top+\Cov\bigl(Z\mathfrak n_t-G_tC\mathfrak n_t\bigr)\succeq\Phi_tP^\perp_t\Phi_t^\top+Q^{\perp\mid y},
\end{equation}
because $Q^\perp-GQ^{y\perp}-Q^{\perp y}G^\top+GQ^yG^\top\succeq Q^{\perp\mid y}$ for every matrix $G$. The order-$d$ model is defined for $t\ge d$: $\hat\xi^{(d)}_t$ is the value at time $t$ of the recursion \cref{eq:apx-perp-filter}, with the exact $\Phi_s$ and $G_s$, started at time $t-d$ from the value $0$. It is a function of $(y_{t-d:t},u_{t-d:t-1})$, and
\begin{equation*}
\hat\xi_t-\hat\xi^{(d)}_t=\Psi_{t,d}\hat\xi_{t-d},\qquad\Psi_{t,d}=\Phi_{t-1}\Phi_{t-2}\cdots\Phi_{t-d},\qquad\Psi_{t,0}=I .
\end{equation*}

\begin{proof}[Proof of \cref{thm:apx-order-d}]
The kernel of the coarse law is $\Ncal(A_{yy}y_t+B_yu_t+A_{y\perp}\hat\xi_t,S_t)$ and that of the order-$d$ model is $\Ncal(A_{yy}y_t+B_yu_t+A_{y\perp}\hat\xi^{(d)}_t,S_t)$, so their expected relative entropy is $\frac12\E^\pi\Vert A_{y\perp}\Psi_{t,d}\hat\xi_{t-d}\Vert^2_{S_t^{-1}}$.

Since $Q\succ0$, $Q^{\perp\mid y}\succ0$ and $\mathsf q_\perp>0$. By \cref{eq:apx-joseph}, $P^\perp_s\succeq\mathsf q_\perp I$ for every $s\ge1$, and for every $s\ge0$, using $P^\perp_{s+1}\preceq\bar pI$,
\begin{equation*}
\Phi_sP^\perp_s\Phi_s^\top\preceq P^\perp_{s+1}-\mathsf q_\perp I\preceq\Bigl(1-\frac{\mathsf q_\perp}{\bar p}\Bigr)P^\perp_{s+1}=\theta P^\perp_{s+1} ,
\end{equation*}
with $\theta\in[0,1)$ because $\bar p\ge\Vert P^\perp_1\Vert\ge\mathsf q_\perp$. Iterating over $s=t-d,\dots,t-1$ gives $\Psi_{t,d}P^\perp_{t-d}\Psi_{t,d}^\top\preceq\theta^dP^\perp_t$. Next, $\E^\pi[\hat\xi_s\hat\xi_s^\top]\preceq\E^\pi[x_{\perp,s}x_{\perp,s}^\top]\preceq\bar X_\perp^2I$, and for $s=t-d\ge1$ the right side is at most $(\bar X_\perp^2/\mathsf q_\perp)P^\perp_{t-d}$. Therefore
\begin{equation*}
\E^\pi\bigl\Vert A_{y\perp}\Psi_{t,d}\hat\xi_{t-d}\bigr\Vert^2_{S_t^{-1}}\le\frac{\bar X_\perp^2}{\mathsf q_\perp}\tr\bigl(S_t^{-1}A_{y\perp}\Psi_{t,d}P^\perp_{t-d}\Psi_{t,d}^\top A_{y\perp}^\top\bigr)\le\theta^d\frac{\bar X_\perp^2}{\mathsf q_\perp}\tr\bigl(S_t^{-1}A_{y\perp}P^\perp_tA_{y\perp}^\top\bigr).
\end{equation*}
If $\lambda_i\ge0$ are the eigenvalues of $(Q^y)^{-1/2}A_{y\perp}P^\perp_tA_{y\perp}^\top(Q^y)^{-1/2}$, the last trace is $\sum_i\lambda_i/(1+\lambda_i)\le\sum_i\ln(1+\lambda_i)=2\mathfrak d_{\mathrm{dyn},t}$, which gives \cref{eq:apx-order-d}.

The restriction $t\ge d+1$ is needed. At $t=d$ the window starts at time $0$, where $P^\perp_0=Z\Pi_C(\Sigma_0)Z^\top$ can be arbitrarily small while $\hat\xi_0$ is not, and the bound then fails (\cref{chk:apx-memory}). If $A_{\perp\perp}=0$ then $\Phi_t=-Q^{\perp y}S_t^{-1}A_{y\perp}$, which is nonzero when $Q^{\perp y}S_t^{-1}A_{y\perp}\ne0$. The predictor of the coarse path can then depend on the entire past although every kernel $\Gamma_l$ with $l\ge2$ vanishes; it does so in the scalar instance of \cref{chk:apx-memory}, where the stationary value of $\Phi_t$ is a nonzero number.
\end{proof}

\begin{lemma}[Comparison for the filter equation]\label{lem:apx-riccati}
Let $\mathsf F$ be a square matrix, $\mathsf J$ a matrix with as many columns, $R\succ0$, and for $P\succeq0$ let $\varphi(P)=P-P\mathsf J^\top(\mathsf JP\mathsf J^\top+R)^{-1}\mathsf JP$ and $\mathcal T_N(P)=\mathsf F\varphi(P)\mathsf F^\top+N$; write $\varphi=\varphi_R$ when the dependence on $R$ matters.
\begin{enumerate}
\item[(a)] $\varphi(P)=\min_K[(I-K\mathsf J)P(I-K\mathsf J)^\top+KRK^\top]$ and $\mathsf F\varphi(P)\mathsf F^\top=\min_K[(\mathsf F-K\mathsf J)P(\mathsf F-K\mathsf J)^\top+KRK^\top]$ in the order of symmetric matrices. Hence $0\preceq\varphi(P)\preceq P$, and $\varphi_R(P)$ is nondecreasing in $P$ and in $R$.
\item[(b)] $\varphi_R(\mu P)=\mu\varphi_{R/\mu}(P)$ for $\mu>0$; hence $\varphi_R(\mu P)\preceq\mu\varphi_R(P)$ for $\mu\ge1$ and $\varphi_R(\mu P)\succeq\mu\varphi_R(P)$ for $0<\mu\le1$.
\item[(c)] If $X_i\succeq0$ solve $X_i=\mathcal T_{N_i}(X_i)$ for $i=1,2$, with $N_2\succ0$ and $N_1\preceq cN_2$ for a constant $c\ge1$, then $X_1\preceq cX_2$. In particular, for $N\succ0$ the equation $X=\mathcal T_N(X)$ has at most one solution $X\succeq0$.
\item[(d)] If $(\mathsf F,\mathsf J)$ is detectable and $N\succ0$, the solution $X\succeq0$ of $X=\mathcal T_N(X)$ exists and is the limit of the nondecreasing iterates $X^{(0)}=0$, $X^{(j+1)}=\mathcal T_N(X^{(j)})$; $X\succeq N$; and with $L=X\mathsf J^\top(\mathsf JX\mathsf J^\top+R)^{-1}=\varphi(X)\mathsf J^\top R^{-1}$, $X=\mathsf F(I-L\mathsf J)X(I-L\mathsf J)^\top\mathsf F^\top+\mathsf FLRL^\top\mathsf F^\top+N$ and the matrices $\mathsf F(I-L\mathsf J)$ and $(I-L\mathsf J)\mathsf F$ have spectral radius less than $1$. If $\rho(\mathsf F)<1$, then $X\preceq Y$ for the solution $Y$ of $Y=\mathsf FY\mathsf F^\top+N$.
\item[(e)] Let $N\succ0$. The following are equivalent: the equation $X=\mathcal T_N(X)$ has a solution $X\succeq0$; the iterates $\mathcal T_N^j(0)$ are bounded; $(\mathsf F,\mathsf J)$ is detectable. When they hold, $\mathcal T_N^j(P_0)\to X$ as $j\to\infty$ for every $P_0\succeq0$.
\end{enumerate}
\end{lemma}

\begin{proof}
(a) With $\Sigma=\mathsf JP\mathsf J^\top+R\succ0$ and $K_\star=P\mathsf J^\top\Sigma^{-1}$, expanding the square gives $(I-K\mathsf J)P(I-K\mathsf J)^\top+KRK^\top=\varphi(P)+(K-K_\star)\Sigma(K-K_\star)^\top$, and in the same way $(\mathsf F-K\mathsf J)P(\mathsf F-K\mathsf J)^\top+KRK^\top=\mathsf F\varphi(P)\mathsf F^\top+(K-\mathsf FK_\star)\Sigma(K-\mathsf FK_\star)^\top$. For fixed $K$ the left sides are nondecreasing in $P$ and in $R$, and a minimum of nondecreasing maps that is attained is nondecreasing. Taking $K=0$ gives $\varphi(P)\preceq P$.

(b) $\varphi_R(\mu P)=\mu P-\mu^2P\mathsf J^\top(\mu\mathsf JP\mathsf J^\top+R)^{-1}\mathsf JP=\mu[P-P\mathsf J^\top(\mathsf JP\mathsf J^\top+R/\mu)^{-1}\mathsf JP]$, and (a) compares $\varphi_{R/\mu}$ with $\varphi_R$.

(c) From $X_2=\mathcal T_{N_2}(X_2)\succeq N_2\succ0$, the set of $\lambda\ge1$ with $X_1\preceq\lambda cX_2$ is nonempty and closed; let $\lambda_\star$ be its least element and suppose $\lambda_\star>1$. By (a) and (b) with $\mu=\lambda_\star c\ge1$,
\begin{equation*}
X_1=\mathcal T_{N_1}(X_1)\preceq\mathsf F\varphi(\lambda_\star cX_2)\mathsf F^\top+N_1\preceq\lambda_\star c\mathsf F\varphi(X_2)\mathsf F^\top+cN_2=\lambda_\star cX_2-(\lambda_\star-1)cN_2 .
\end{equation*}
With $\nu=\lambda_{\min}(N_2)/\Vert X_2\Vert\in(0,1]$ one has $N_2\succeq\nu X_2$, so $X_1\preceq(\lambda_\star-(\lambda_\star-1)\nu)cX_2$ with $1\le\lambda_\star-(\lambda_\star-1)\nu<\lambda_\star$, which contradicts the choice of $\lambda_\star$. Hence $\lambda_\star=1$. Uniqueness follows by applying this twice with $N_1=N_2=N$ and $c=1$.

(d) Let $K_0$ be such that $\rho(\mathsf F-K_0\mathsf J)<1$ and let $Y_0\succeq0$ solve $Y_0=(\mathsf F-K_0\mathsf J)Y_0(\mathsf F-K_0\mathsf J)^\top+K_0RK_0^\top+N$. The iterates $X^{(0)}=0$, $X^{(j+1)}=\mathcal T_N(X^{(j)})$ are nondecreasing, because $X^{(1)}=N\succeq X^{(0)}$ and $\mathcal T_N$ is nondecreasing, and they satisfy $X^{(j)}\preceq Y_0$ by induction, because (a) gives $\mathcal T_N(P)\preceq(\mathsf F-K_0\mathsf J)P(\mathsf F-K_0\mathsf J)^\top+K_0RK_0^\top+N$. They converge, and the limit solves $X=\mathcal T_N(X)$ by continuity of $\varphi$. If $\rho(\mathsf F)<1$ one may take $K_0=0$, which gives $X\preceq Y$. By (a) with $K=\mathsf FL$, $X=\mathsf F(I-L\mathsf J)X(I-L\mathsf J)^\top\mathsf F^\top+\mathsf FLRL^\top\mathsf F^\top+N$. If $v^*\mathsf F(I-L\mathsf J)=\lambda v^*$ with $v\ne0$, then $(1-\lvert\lambda\rvert^2)v^*Xv\ge v^*Nv>0$, so $\lvert\lambda\rvert<1$; and $(I-L\mathsf J)\mathsf F$ has the same nonzero eigenvalues. Finally $\varphi(X)\mathsf J^\top=X\mathsf J^\top\Sigma^{-1}(\Sigma-\mathsf JX\mathsf J^\top)=X\mathsf J^\top\Sigma^{-1}R$ with $\Sigma=\mathsf JX\mathsf J^\top+R$.

(e) If $(\mathsf F,\mathsf J)$ is detectable, a solution exists by (d). If $X\succeq0$ is a solution, then $\mathcal T_N^j(0)\preceq\mathcal T_N^j(X)=X$ for every $j$, because $\mathcal T_N$ is nondecreasing, so the iterates are bounded. If the iterates are bounded, they converge, being nondecreasing, and the limit $X$ solves $X=\mathcal T_N(X)$ by continuity of $\varphi$; the argument of (d) for the spectral radius uses only this equation, the identity of (a) with $K=\mathsf FL$, and $X\succeq N\succ0$, so $\rho(\mathsf F-\mathsf FL\mathsf J)<1$, which is detectability. For the convergence, let $X$ be the solution and $P_0\succeq0$. Since $X\succeq N\succ0$ there is $\mu\ge1$ with $P_0\preceq\mu X$, and by (a) and (b) $\mathcal T_N(\mu X)\preceq\mu\mathsf F\varphi(X)\mathsf F^\top+N\preceq\mu X$. Hence the iterates $\mathcal T_N^j(\mu X)$ are nonincreasing; they are bounded below by $\mathcal T_N^j(X)=X$, so they converge to a solution $U\succeq0$ of $U=\mathcal T_N(U)$, which is $X$ by (c). The iterates $\mathcal T_N^j(0)$ increase to a solution, which is again $X$. Since $\mathcal T_N^j(0)\preceq\mathcal T_N^j(P_0)\preceq\mathcal T_N^j(\mu X)$, $\mathcal T_N^j(P_0)\to X$.
\end{proof}

The fine filter of \cref{lem:apx-separation}(b) is the case $\mathsf F=A$, $\mathsf J=HC$, $N=Q$, and the declared filter is the case $\mathsf F=\bar A$, $\mathsf J=H$, $N=\bar Q$, for which $\varphi=\varphi_R$ of the main text. Since $\rho(A)<1$, \cref{lem:apx-riccati}(d) gives a unique $\Sigma^-\succeq0$ with $\Sigma^-=A\Sigma^{\mathrm f}A^\top+Q$, where $\Sigma^{\mathrm f}$ is obtained from $\Sigma^-$ by the middle formula of \cref{eq:apx-record-filter}, and $\Sigma^{\mathrm f}\preceq\Sigma^-\preceq\Sigma_x$. That the fine filter is stationary means $\Sigma_0=\Sigma^-$, so that $\Sigma^-_t=\Sigma^-$ and $\Sigma^{\mathrm f}_t=\Sigma^{\mathrm f}$ for all $t$. Detectability of $(\bar A,H)$ gives $P^-_{\mathrm b}$, with $P^-_{\mathrm b}\succeq\bar Q\succ0$, $P_{\mathrm b}=\varphi_R(P^-_{\mathrm b})=((P^-_{\mathrm b})^{-1}+H^\top R^{-1}H)^{-1}\succ0$, $L=P^-_{\mathrm b}H^\top(HP^-_{\mathrm b}H^\top+R)^{-1}=P_{\mathrm b}H^\top R^{-1}$, and $\rho(\Psi')=\rho(\Psi)<1$ for $\Psi'=\bar A(I-LH)$ and $\Psi=(I-LH)\bar A$. The declared filter is $\hat\eta^-_{t+1}=\bar A(\hat\eta^-_t+L(s_{t+1}-H\hat\eta^-_t))+\bar Bu_t$, started from any $\hat\eta^-_0$ with $\E\lvert C\hat x^-_0-\hat\eta^-_0\rvert^2<\infty$.

\begin{proof}[Proof of \cref{thm:apx-posterior}]
Fixed-point identity. By \cref{lem:apx-separation}(b) in the stationary regime, $P_{\mathrm e}=C\Sigma^-C^\top=CA\Sigma^{\mathrm f}A^\top C^\top+Q^y$. Inserting $CA=\bar AC+D$ gives $CA\Sigma^{\mathrm f}A^\top C^\top=\bar AC\Sigma^{\mathrm f}C^\top\bar A^\top+\Delta$, and $C\Sigma^{\mathrm f}C^\top=P_{\mathrm e}-P_{\mathrm e}H^\top(HP_{\mathrm e}H^\top+R)^{-1}HP_{\mathrm e}=\varphi_R(P_{\mathrm e})$. Hence
\begin{equation}\label{eq:apx-fixed-point}
P_{\mathrm e}=\bar A\varphi_R(P_{\mathrm e})\bar A^\top+Q^y+\Delta ,
\end{equation}
the filter equation of the declared model with $\bar Q$ replaced by $Q^y+\Delta$, which is symmetric and satisfies $Q^y+\Delta\succeq c_-\bar Q\succ0$.

Sandwich. \Cref{lem:apx-riccati}(c) with $X_1=P_{\mathrm e}$, $N_1=Q^y+\Delta$, $X_2=P^-_{\mathrm b}$, $N_2=\bar Q$ and $c=c_+$ gives $P_{\mathrm e}\preceq c_+P^-_{\mathrm b}$. With $X_1=P^-_{\mathrm b}$, $N_1=\bar Q$, $X_2=P_{\mathrm e}$, $N_2=Q^y+\Delta$ and $c=1/c_-\ge1$ it gives $P^-_{\mathrm b}\preceq P_{\mathrm e}/c_-$. This step needs $c_->0$.

Decomposition. Let $e_t=C\hat x^-_t-\hat\eta^-_t$. Both laws in $\Delta\mathcal F_t$ are Gaussian, so
\begin{equation}\label{eq:apx-excess-split}
\Delta\mathcal F_t=h\bigl((P^-_{\mathrm b})^{-1}P_{\mathrm e}\bigr)+\tfrac12\E\Vert e_t\Vert^2_{(P^-_{\mathrm b})^{-1}} .
\end{equation}
The eigenvalues of $(P^-_{\mathrm b})^{-1}P_{\mathrm e}$ lie in $[c_-,c_+]$ and $g$ is decreasing on $(0,1]$ and increasing on $[1,\infty)$, so the first term is at most $\frac k2\max\{g(c_+),g(c_-)\}$.

Error recursion. The innovation $\iota_t=s_{t+1}-HC\hat x^-_t=HC(x_t-\hat x^-_t)+v_t$ is, given $(s_{1:t},\omega)$, Gaussian with mean zero and covariance $\Sigma_\iota=HP_{\mathrm e}H^\top+R$; the $\iota_t$ are therefore independent and identically distributed. With $L_{\mathrm e}=P_{\mathrm e}H^\top\Sigma_\iota^{-1}$, $C\hat x^{\mathrm f}_t=C\hat x^-_t+L_{\mathrm e}\iota_t$, and $C\hat x^-_{t+1}=CA\hat x^{\mathrm f}_t+CBu_t=\bar AC\hat x^{\mathrm f}_t+\bar Bu_t+d_t$ with $d_t=D\hat x^{\mathrm f}_t+Eu_t$. Since $s_{t+1}-H\hat\eta^-_t=\iota_t+He_t$, subtracting the declared filter gives
\begin{equation}\label{eq:apx-error-recursion}
e_{t+1}=\Psi'e_t+\bar A(L_{\mathrm e}-L)\iota_t+d_t .
\end{equation}
The act $u_t$ is measurable in $(s_{1:t},\omega)$, so $d_t=\E[\delta_t\mid s_{1:t+1},\omega]$ and $\Vert d_t\Vert_{L^2}\le\Vert\delta_t\Vert_{L^2}$.

Gain difference. By \cref{lem:apx-riccati}(d) applied to \cref{eq:apx-fixed-point}, $L_{\mathrm e}=\varphi_R(P_{\mathrm e})H^\top R^{-1}$, so $L_{\mathrm e}-L=WH^\top R^{-1}$ with $W=\varphi_R(P_{\mathrm e})-\varphi_R(P^-_{\mathrm b})$. By the sandwich and \cref{lem:apx-riccati}(a),(b), $c_-P_{\mathrm b}\preceq\varphi_R(P_{\mathrm e})\preceq c_+P_{\mathrm b}$, hence $-\mathfrak cP_{\mathrm b}\preceq W\preceq\mathfrak cP_{\mathrm b}$ and $\Vert W\Vert\le\mathfrak c\Vert P_{\mathrm b}\Vert$. Since $\Sigma_\iota\preceq c_+HP^-_{\mathrm b}H^\top+R$,
\begin{equation*}
\tr\bigl((L_{\mathrm e}-L)\Sigma_\iota(L_{\mathrm e}-L)^\top\bigr)=\tr\bigl(WH^\top R^{-1}\Sigma_\iota R^{-1}HW\bigr)\le k\mathfrak c^2\Vert P_{\mathrm b}\Vert^2\mathfrak g^2 .
\end{equation*}

Conclusion. Solving \cref{eq:apx-error-recursion} and using the independence of the innovations,
\begin{equation*}
\Vert e_t\Vert_{L^2}\le\Vert\Psi'^t\Vert\Vert e_0\Vert_{L^2}+\sum_{s<t}\Vert\Psi'^{t-1-s}\Vert\Vert\delta_s\Vert_{L^2}+\Bigl(\sum_{j<t}\Vert\Psi'^j\bar A\Vert^2\Bigr)^{1/2}\sqrt k\mathfrak c\Vert P_{\mathrm b}\Vert\mathfrak g .
\end{equation*}
Since $\rho(\Psi')<1$, the sums $\mathcal S'$ and $\mathcal S'_2$ are finite, the first term tends to zero, and
\begin{equation*}
\limsup_t\sum_{s<t}\Vert\Psi'^{t-1-s}\Vert\Vert\delta_s\Vert_{L^2}\le\mathcal S'\limsup_s\Vert\delta_s\Vert_{L^2}\le\mathcal S'\bar{\mathfrak d}.
\end{equation*}
With $\Vert e_t\Vert^2_{(P^-_{\mathrm b})^{-1}}\le\lvert e_t\rvert^2/\lambda_{\min}(P^-_{\mathrm b})$ and \cref{eq:apx-excess-split} this is \cref{eq:apx-posterior}.

A priori constants. Since $\Sigma^{\mathrm f}\preceq\Sigma_x$, $\Vert D\Sigma^{\mathrm f}D^\top\Vert\le\sigma_D^2$ and $\Vert C\Sigma^{\mathrm f}D^\top\Vert\le\Vert C\Sigma^{\mathrm f}C^\top\Vert^{1/2}\Vert D\Sigma^{\mathrm f}D^\top\Vert^{1/2}\le\varsigma_y\sigma_D$, so $\Vert\Delta\Vert\le\vartheta$ and $Q^y-\vartheta I\preceq Q^y+\Delta\preceq Q^y+\vartheta I$. With $\vartheta I\preceq(\vartheta/\lambda_{\min}(\bar Q))\bar Q$ the constants of the statement satisfy $c_-\bar Q\preceq Q^y+\Delta\preceq c_+\bar Q$; they are admissible when the resulting $c_-$ is positive, and the theorem gives nothing otherwise.
\end{proof}

\begin{proposition}[Filtered law]\label{prop:apx-posterior-filtered}
Under the hypotheses of \cref{thm:apx-posterior}, the exact law of $y_t$ given $(s_{1:t+1},\omega)$ is $\Ncal(C\hat x^{\mathrm f}_t,\varphi_R(P_{\mathrm e}))$, the declared filter holds $\Ncal(\hat\eta^{\mathrm f}_t,P_{\mathrm b})$ with $\hat\eta^{\mathrm f}_t=\hat\eta^-_t+L(s_{t+1}-H\hat\eta^-_t)$, one has $c_-P_{\mathrm b}\preceq\varphi_R(P_{\mathrm e})\preceq c_+P_{\mathrm b}$, and
\begin{equation}\label{eq:apx-posterior-filtered}
\limsup_{t\to\infty}\E\KL\Bigl(\Ncal\bigl(C\hat x^{\mathrm f}_t,\varphi_R(P_{\mathrm e})\bigr)\Big\Vert\Ncal\bigl(\hat\eta^{\mathrm f}_t,P_{\mathrm b}\bigr)\Bigr)\le\frac k2\max\{g(c_+),g(c_-)\}+\frac{\bigl(\mathcal S^{\mathrm f}\bar{\mathfrak d}+\mathfrak c\Vert P_{\mathrm b}\Vert\mathfrak g\sqrt k\mathcal S^{\mathrm f}_2\bigr)^2}{2\lambda_{\min}(P_{\mathrm b})},
\end{equation}
with $\mathcal S^{\mathrm f}=\sum_{j\ge0}\Vert\Psi^j(I-LH)\Vert$ and $\mathcal S^{\mathrm f}_2=(\sum_{j\ge0}\Vert\Psi^j\Vert^2)^{1/2}$.
\end{proposition}

\begin{proof}
The covariance statement is the sandwich for $\varphi_R(P_{\mathrm e})$ from the previous proof. Let $e^{\mathrm f}_t=C\hat x^{\mathrm f}_t-\hat\eta^{\mathrm f}_t=(I-LH)e_t+(L_{\mathrm e}-L)\iota_t$. By \cref{eq:apx-error-recursion},
\begin{equation*}
e^{\mathrm f}_{t+1}=\Psi e^{\mathrm f}_t+(I-LH)d_t+(L_{\mathrm e}-L)\iota_{t+1},
\end{equation*}
so $e^{\mathrm f}_t=\Psi^t(I-LH)e_0+\sum_{s<t}\Psi^{t-1-s}(I-LH)d_s+\sum_{s\le t}\Psi^{t-s}(L_{\mathrm e}-L)\iota_s$, and the three terms are bounded as before, with $\lambda_{\min}(P_{\mathrm b})$ in place of $\lambda_{\min}(P^-_{\mathrm b})$ because the reference covariance is now $P_{\mathrm b}$.
\end{proof}

The constants of \cref{eq:apx-posterior-filtered} differ from those of \cref{eq:apx-posterior}: $\lambda_{\min}(P_{\mathrm b})\le\lambda_{\min}(P^-_{\mathrm b})$, and the factor $I-LH$ cannot be removed from $\mathcal S^{\mathrm f}$ at no cost, since its spectral norm can exceed $1$. Whether \cref{eq:apx-posterior} holds verbatim for the filtered law with $\Psi$ in place of $\Psi'$ was neither proved nor refuted.

\begin{proof}[Proof of \cref{cor:apx-memory}]
Without input, $B=0$, $\bar B=0$ and $\delta_t=Dx_t$. Write $\Sigma_y=C\Sigma_xC^\top$ and $\Sigma_{x\mid y}=\Pi_C(\Sigma_x)$. The moment-matched $\bar A=CA\Sigma_xC^\top\Sigma_y^{-1}$ gives $D\Sigma_xC^\top=CA\Sigma_xC^\top-\bar A\Sigma_y=0$, and therefore
\begin{equation*}
D\Sigma_xD^\top=CA\Sigma_xA^\top C^\top-CA\Sigma_xC^\top\Sigma_y^{-1}C\Sigma_xA^\top C^\top=CA\Sigma_{x\mid y}A^\top C^\top .
\end{equation*}
In the stationary law, $x_t$ given $y_t$ is Gaussian with covariance $\Sigma_{x\mid y}$ and $y_{t+1}=CAx_t+C\mathfrak n_t$, so the closure defect $\mathfrak d_{\mathrm{cl}}=I(y_{t+1};x_t\mid y_t)$ of \cref{def:agy-closure}(i), here without a blanket variable, is
\begin{equation*}
\mathfrak d_{\mathrm{cl}}=\tfrac12\ln\det\bigl(I+(Q^y)^{-1}D\Sigma_xD^\top\bigr),
\end{equation*}
the Gaussian form of \cref{prop:met-closure-notions}(a) and \cref{prop:met-memory}(b). The largest eigenvalue of $(Q^y)^{-1/2}D\Sigma_xD^\top(Q^y)^{-1/2}$ is at least $\sigma_D^2/\Vert Q^y\Vert$, and the determinant is at least $1$ plus that eigenvalue, so $e^{2\mathfrak d_{\mathrm{cl}}}\ge1+\sigma_D^2/\Vert Q^y\Vert$, which is the first claim.

Defect energy. The fine filter is stationary, so $\Cov(x_0)=\Sigma^-\preceq\Sigma_x$, and by induction $\Cov(x_t)\preceq\Sigma_x$ for all $t$, because $A\Sigma_xA^\top+Q=\Sigma_x$. The mean of $x_t$ is $A^tm_0$, which tends to zero. Hence $\limsup_t\E\lvert Dx_t\rvert^2\le\tr(D\Sigma_xD^\top)\le k\sigma_D^2$, and $\bar{\mathfrak d}=\sqrt k\sigma_D$ is admissible in \cref{thm:apx-posterior}. If $m_0=0$ the bound holds for every $t$.

Detectability. Since $D\Sigma_xC^\top=0$, $\Sigma_y=CA\Sigma_xA^\top C^\top+Q^y=\bar A\Sigma_y\bar A^\top+D\Sigma_xD^\top+Q^y$ with $Q^y\succ0$, so $\rho(\bar A)<1$ and $(\bar A,H)$ is detectable for every $H$.

Constants. With $\bar Q=Q^y$ the a priori constants of \cref{thm:apx-posterior} are $c_\pm=1\pm\mathsf x$, which are admissible when $\mathsf x<1$, and $\mathfrak c=\mathsf x$. From the power series, $g(1+\mathsf x)=\mathsf x-\ln(1+\mathsf x)\le\mathsf x^2/2$ and $g(1-\mathsf x)=\sum_{j\ge2}\mathsf x^j/j\le\mathsf x^2/(2(1-\mathsf x))$, so $\frac k2\max\{g(c_+),g(c_-)\}\le\frac k4\mathsf x^2/(1-\mathsf x)$. By \cref{lem:apx-riccati}(d), $\lambda_{\min}(P^-_{\mathrm b})\ge\lambda_{\min}(Q^y)$. Inserting these in \cref{eq:apx-posterior} gives the stated bound.

Order. By the first claim $\sigma_D^2=O(\mathfrak d_{\mathrm{cl}})$, hence $\vartheta=O(\mathfrak d_{\mathrm{cl}}^{1/2})$ and $\mathsf x^2=O(\mathfrak d_{\mathrm{cl}})$, and the bound is $O(\mathfrak d_{\mathrm{cl}})$ along every family of instances on which $\Vert\bar A\Vert$, $\varsigma_y$, $\Vert Q^y\Vert$, $\lambda_{\min}(Q^y)^{-1}$, $\mathcal S'$, $\mathcal S'_2$, $\Vert P_{\mathrm b}\Vert$ and $\mathfrak g$ stay bounded. Without such a restriction the constants can grow as the defect shrinks, and no order is claimed.
\end{proof}

\begin{lemma}[Difference of filtered covariances]\label{lem:apx-difference}
In the notation of \cref{lem:apx-riccati} with $\mathsf J=H$, let $P_1,P_2\succeq0$, $\mathsf S_i=HP_iH^\top+R$, $L_i=P_iH^\top\mathsf S_i^{-1}$, $J_i=H^\top\mathsf S_i^{-1}H$ and $X=P_1-P_2$. Then
\begin{equation}\label{eq:apx-difference}
\begin{gathered}
\varphi(P_1)-\varphi(P_2)=(I-L_2H)\bigl(X-XJ_1X\bigr)(I-L_2H)^\top=(I-L_1H)X(I-L_2H)^\top,\\
L_1-L_2=(I-L_2H)XH^\top\mathsf S_1^{-1}.
\end{gathered}
\end{equation}
\end{lemma}

\begin{proof}
Gain difference. Since $(I-L_2H)P_2H^\top=P_2H^\top\mathsf S_2^{-1}(\mathsf S_2-HP_2H^\top)=L_2R$ and $L_2HP_1H^\top\mathsf S_1^{-1}=L_2(\mathsf S_1-R)\mathsf S_1^{-1}=L_2-L_2R\mathsf S_1^{-1}$,
\begin{equation*}
(I-L_2H)XH^\top\mathsf S_1^{-1}=L_1-L_2HP_1H^\top\mathsf S_1^{-1}-L_2R\mathsf S_1^{-1}=L_1-L_2 .
\end{equation*}
First form. The identity in the proof of \cref{lem:apx-riccati}(a), $(I-KH)P(I-KH)^\top+KRK^\top=\varphi(P)+(K-K_\star)(HPH^\top+R)(K-K_\star)^\top$ with $K_\star=PH^\top(HPH^\top+R)^{-1}$, applied with $K=L_2$ to $P=P_1$ and to $P=P_2$, gives $\varphi(P_1)=(I-L_2H)P_1(I-L_2H)^\top+L_2RL_2^\top-(L_2-L_1)\mathsf S_1(L_2-L_1)^\top$ and $\varphi(P_2)=(I-L_2H)P_2(I-L_2H)^\top+L_2RL_2^\top$. Subtracting and inserting the gain difference, $(L_1-L_2)\mathsf S_1(L_1-L_2)^\top=(I-L_2H)XJ_1X(I-L_2H)^\top$, which is the first form. Second form. $\mathsf S_1^{-1}-\mathsf S_2^{-1}=\mathsf S_1^{-1}(\mathsf S_2-\mathsf S_1)\mathsf S_2^{-1}=-\mathsf S_1^{-1}HXH^\top\mathsf S_2^{-1}$, so $J_1-J_2=-J_1XJ_2$, and therefore
\begin{equation*}
(I-J_1X)(I-J_2P_2)=I-J_1P_1+\bigl(J_1-J_2+J_1XJ_2\bigr)P_2=I-J_1P_1 .
\end{equation*}
Since $(I-L_iH)^\top=I-J_iP_i$, this reads $(I-J_1X)(I-L_2H)^\top=(I-L_1H)^\top$, so the first form equals $(I-L_2H)X(I-L_1H)^\top$, and the left side of \cref{eq:apx-difference} is symmetric, so it also equals the transpose $(I-L_1H)X(I-L_2H)^\top$.
\end{proof}

\begin{proof}[Proof of \cref{thm:apx-unconditional}]
Write $P^-=P^-_{\mathrm b}$ and $\mathcal T(P)=\bar A\varphi_R(P)\bar A^\top+\bar Q$, so that $\mathcal T(P^-)=P^-$. The fixed-point identity \cref{eq:apx-fixed-point} and the error recursion \cref{eq:apx-error-recursion} were derived without the hypothesis $c_->0$, and $P_{\mathrm e}=C\Sigma^-C^\top\succeq Q^y\succ0$.

(a) By \cref{lem:apx-riccati}(d), $P^-=\Psi'P^-\Psi'^\top+\bar ALRL^\top\bar A^\top+\bar Q$, so $\Psi'P^-\Psi'^\top\preceq P^--\bar Q\preceq\theta^0_{\mathrm b}P^-$, because $\bar Q\succeq\lambda_{\min}(\tilde{\bar Q})P^-=(1-\theta^0_{\mathrm b})P^-$. With $\tilde\Psi=(P^-)^{-1/2}\Psi'(P^-)^{1/2}$ this says $\tilde\Psi\tilde\Psi^\top\preceq\theta^0_{\mathrm b}I$, hence $\theta_{\mathrm b}=\Vert\tilde\Psi\Vert^2\le\theta^0_{\mathrm b}$, and $\theta^0_{\mathrm b}<1$ because $\bar Q\succ0$. Consequently, for every symmetric $M\succeq0$ with $M\preceq\mathsf mP^-$ and every $j\ge0$,
\begin{equation}\label{eq:apx-metric-contraction}
\Psi'^jM\Psi'^{j\top}\preceq\mathsf m\theta_{\mathrm b}^jP^-,\qquad\tr\bigl((P^-)^{-1}\Psi'^jM\Psi'^{j\top}\bigr)\le\theta_{\mathrm b}^j\tr\tilde M ,
\end{equation}
since $(P^-)^{-1/2}\Psi'^jM\Psi'^{j\top}(P^-)^{-1/2}=\tilde\Psi^j\tilde M\tilde\Psi^{j\top}$. \Cref{lem:apx-difference} with $P_1=P_{\mathrm e}$, $P_2=P^-$, for which $L_1=L_{\mathrm e}$, $L_2=L$, $\mathsf S_1=\Sigma_\iota$ and $J_1=J_{\mathrm e}$, gives $\bar A\varphi_R(P_{\mathrm e})\bar A^\top-\bar A\varphi_R(P^-)\bar A^\top=\Psi'(\mathsf X-\mathsf XJ_{\mathrm e}\mathsf X)\Psi'^\top$. Subtracting $P^-=\bar A\varphi_R(P^-)\bar A^\top+\bar Q$ from \cref{eq:apx-fixed-point} gives \cref{eq:apx-quadratic}. The same lemma gives $\bar A(L_{\mathrm e}-L)=\Psi'\mathsf XH^\top\Sigma_\iota^{-1}$, and $\iota_t$ has covariance $\Sigma_\iota$, so $\bar A(L_{\mathrm e}-L)\iota_t$ has covariance $\Psi'\mathsf XJ_{\mathrm e}\mathsf X\Psi'^\top$.

(b) Upper side and second-order term. By \cref{eq:apx-quadratic}, $\mathsf X$ solves $\mathsf X=\Psi'\mathsf X\Psi'^\top+\mathsf N-\Psi'\mathsf XJ_{\mathrm e}\mathsf X\Psi'^\top$, and $\rho(\Psi')<1$, so $\mathsf X=\mathsf X_1-\sum_{j\ge1}\Psi'^j\mathsf XJ_{\mathrm e}\mathsf X\Psi'^{j\top}$. The sum is positive semidefinite, which gives $\mathsf X\preceq\mathsf X_1$ and $\mathsf Y\preceq\mathsf Y_1$. Since $\Sigma_\iota\succeq R$, $(P^-)^{1/2}J_{\mathrm e}(P^-)^{1/2}\preceq\mathsf jI$, so $\mathsf XJ_{\mathrm e}\mathsf X=(P^-)^{1/2}\mathsf Y(P^-)^{1/2}J_{\mathrm e}(P^-)^{1/2}\mathsf Y(P^-)^{1/2}\preceq\mathsf j\Vert\mathsf Y\Vert^2P^-$, and \cref{eq:apx-metric-contraction} gives $\mathsf X_1-\mathsf X\preceq\mathsf j\Vert\mathsf Y\Vert^2\sum_{j\ge1}\theta_{\mathrm b}^jP^-$, the left inequality of \cref{eq:apx-enclosure}. For $\mathsf Y_1=\sum_j\tilde\Psi^j\tilde{\mathsf N}\tilde\Psi^{j\top}$, $\Vert\mathsf Y_1\Vert\le\Vert\tilde{\mathsf N}\Vert\sum_j\theta_{\mathrm b}^j$. From $P_{\mathrm e}\succeq Q^y\succeq\mathsf lP^-$, $\mathsf Y=\tilde P_{\mathrm e}-I\succeq(\mathsf l-1)I$.

Norm bound. Let $\mathsf P_0=P_{\mathrm e}$ and $\mathsf P_{i+1}=\mathcal T(\mathsf P_i)$. By \cref{eq:apx-fixed-point}, $\mathsf P_1=P_{\mathrm e}-\mathsf N$, and $\mathsf P_1=\mathcal T(P_{\mathrm e})\succeq\bar Q\succ0$. Since $\mathsf Y\preceq\mathsf Y_1$, $P_{\mathrm e}\preceq\mathsf uP^-$ with $\mathsf u\ge1$, and by \cref{lem:apx-riccati}(a),(b) $\mathcal T(\mathsf uP^-)\preceq\mathsf u\bar A\varphi_R(P^-)\bar A^\top+\bar Q\preceq\mathsf uP^-$; $\mathcal T$ is nondecreasing, so by induction $\mathsf P_i\preceq\mathcal T^i(\mathsf uP^-)\preceq\mathsf uP^-$ for all $i$. The iterates $\mathcal T^i(\mathsf uP^-)$ are nonincreasing and bounded below by $\mathcal T^i(P^-)=P^-$, so they converge to a solution of $U=\mathcal T(U)$ with $U\succeq0$, which is $P^-$ by \cref{lem:apx-riccati}(c); the iterates $\mathcal T^i(0)$ increase to $P^-$ by \cref{lem:apx-riccati}(d); and $\mathcal T^i(0)\preceq\mathsf P_i\preceq\mathcal T^i(\mathsf uP^-)$. Hence $\mathsf P_i\to P^-$ and $P^--P_{\mathrm e}=\sum_{i\ge0}(\mathsf P_{i+1}-\mathsf P_i)$. Let $\Psi_i=\bar A(I-L(\mathsf P_i)H)$ with $L(P)=PH^\top(HPH^\top+R)^{-1}$. By the second form of \cref{eq:apx-difference}, $\mathsf P_{i+1}-\mathsf P_i=\Psi_i(\mathsf P_i-\mathsf P_{i-1})\Psi_{i-1}^\top$ for $i\ge1$, so
\begin{equation*}
\mathsf P_{i+1}-\mathsf P_i=-\mathsf A_i\mathsf N\mathsf B_i^\top,\qquad\mathsf A_i=\Psi_i\Psi_{i-1}\cdots\Psi_1,\quad\mathsf B_i=\Psi_{i-1}\cdots\Psi_0,\quad\mathsf A_0=\mathsf B_0=I .
\end{equation*}
By \cref{lem:apx-riccati}(a) with $K=\bar AL(\mathsf P_i)$, $\mathsf P_{i+1}=\Psi_i\mathsf P_i\Psi_i^\top+\bar AL(\mathsf P_i)RL(\mathsf P_i)^\top\bar A^\top+\bar Q\succeq\Psi_i\mathsf P_i\Psi_i^\top+\bar Q$, and $\bar Q\succeq(1-\theta^0_{\mathrm b})P^-\succeq\mathsf u^{-1}(1-\theta^0_{\mathrm b})\mathsf P_{i+1}$. Hence $\Psi_i\mathsf P_i\Psi_i^\top\preceq\chi_{\mathrm b}\mathsf P_{i+1}$ with $\chi_{\mathrm b}=1-(1-\theta^0_{\mathrm b})/\mathsf u\in[0,1)$, and by induction $\mathsf B_i\mathsf P_0\mathsf B_i^\top\preceq\chi_{\mathrm b}^i\mathsf P_i$ and $\mathsf A_i\mathsf P_1\mathsf A_i^\top\preceq\chi_{\mathrm b}^i\mathsf P_{i+1}$. For every vector $z$, with $\mathsf w=\Vert\mathsf P_1^{-1/2}\mathsf N\mathsf P_0^{-1/2}\Vert$,
\begin{equation*}
\bigl\lvert z^\top\mathsf A_i\mathsf N\mathsf B_i^\top z\bigr\rvert\le\mathsf w\bigl(z^\top\mathsf A_i\mathsf P_1\mathsf A_i^\top z\bigr)^{1/2}\bigl(z^\top\mathsf B_i\mathsf P_0\mathsf B_i^\top z\bigr)^{1/2}\le\mathsf w\chi_{\mathrm b}^i\bigl(z^\top\mathsf P_{i+1}z\bigr)^{1/2}\bigl(z^\top\mathsf P_iz\bigr)^{1/2}\le\mathsf w\mathsf u\chi_{\mathrm b}^iz^\top P^-z .
\end{equation*}
Summing over $i$, $\lvert z^\top\mathsf Xz\rvert\le\mathsf w\mathsf u(1-\chi_{\mathrm b})^{-1}z^\top P^-z=\mathsf rz^\top P^-z$, that is $\Vert\mathsf Y\Vert\le\mathsf r$. Finally $\mathsf P_1\succeq\bar Q\succeq(1-\theta^0_{\mathrm b})P^-$ and $\mathsf P_0\succeq\mathsf lP^-$ give $\Vert\mathsf P_1^{-1/2}(P^-)^{1/2}\Vert\le(1-\theta^0_{\mathrm b})^{-1/2}$ and $\Vert(P^-)^{1/2}\mathsf P_0^{-1/2}\Vert\le\mathsf l^{-1/2}$, hence $\mathsf w\le\Vert\tilde{\mathsf N}\Vert/((1-\theta^0_{\mathrm b})\mathsf l)^{1/2}$.

(c) By \cref{eq:apx-excess-split}, $\Delta\mathcal F_t=h((P^-)^{-1}P_{\mathrm e})+\frac12\E\Vert e_t\Vert^2_{(P^-)^{-1}}$, and the eigenvalues of $(P^-)^{-1}P_{\mathrm e}$ are $1+\mathsf y_i$, so the first term is $\frac12\sum_ig(1+\mathsf y_i)$. For $\mathsf y\ge0$, $g(1+\mathsf y)\le\mathsf y^2/2$, and for $-1<\mathsf y<0$, $g(1+\mathsf y)=\sum_{j\ge2}(-\mathsf y)^j/j\le\mathsf y^2/(2(1+\mathsf y))$, which gives the stated bound on the first term; $1+\lambda_{\min}(\mathsf Y)\ge\max\{\mathsf l,1-\mathsf r\}$ and $\Vert\mathsf Y\Vert_{\mathrm F}\le\sqrt k\Vert\mathsf Y\Vert$ by (b). For the second term write $\Vert v\Vert_{\mathrm b}=(\E\Vert v\Vert^2_{(P^-)^{-1}})^{1/2}$. Solving \cref{eq:apx-error-recursion},
\begin{equation*}
e_t=\Psi'^te_0+\sum_{s<t}\Psi'^{t-1-s}d_s+\sum_{s<t}\Psi'^{t-1-s}\bar A(L_{\mathrm e}-L)\iota_s .
\end{equation*}
By \cref{eq:apx-metric-contraction} in the form $\Vert(P^-)^{-1/2}\Psi'^jv\Vert\le\theta_{\mathrm b}^{j/2}\Vert(P^-)^{-1/2}v\Vert$, the first term tends to zero and, since $\E\lvert u_s\rvert^2<\infty$ makes every $\Vert d_s\Vert_{\mathrm b}$ finite, the second has $\Vert\cdot\Vert_{\mathrm b}$ at most $\sum_{s<t}\theta_{\mathrm b}^{(t-1-s)/2}\Vert d_s\Vert_{\mathrm b}$, whose upper limit is at most $\tau_{\mathrm b}\bar d_{\mathrm b}$. The innovations are independent with mean zero, so by (a) and \cref{eq:apx-metric-contraction} the third term has
\begin{equation*}
\begin{aligned}
\Bigl\Vert\sum_{s<t}\Psi'^{t-1-s}\bar A(L_{\mathrm e}-L)\iota_s\Bigr\Vert_{\mathrm b}^2&=\sum_{j<t}\tr\bigl((P^-)^{-1}\Psi'^{j+1}\mathsf XJ_{\mathrm e}\mathsf X\Psi'^{(j+1)\top}\bigr)\\
&\le\sum_{j\ge0}\theta_{\mathrm b}^{j+1}\tr\bigl(\mathsf Y(P^-)^{1/2}J_{\mathrm e}(P^-)^{1/2}\mathsf Y\bigr)\le\frac{\theta_{\mathrm b}\mathsf j}{1-\theta_{\mathrm b}}\Vert\mathsf Y\Vert_{\mathrm F}^2 .
\end{aligned}
\end{equation*}
The triangle inequality for $\Vert\cdot\Vert_{\mathrm b}$ gives \cref{eq:apx-unconditional}. Since $u_t$ is measurable in $(s_{1:t},\omega)$, $d_t=\E[\delta_t\mid s_{1:t+1},\omega]$, a conditional expectation, which does not increase $\Vert\cdot\Vert_{\mathrm b}$, and $(P^-)^{-1}\preceq\bar Q^{-1}$, so $\E\Vert d_t\Vert^2_{(P^-)^{-1}}\le\E\Vert\delta_t\Vert^2_{\bar Q^{-1}}$. Without input $d_t=D\hat x^{\mathrm f}_t$ with $\hat x^{\mathrm f}_t=\E[x_t\mid s_{1:t+1},\omega]$, and $\Cov(x_t)=\Cov(\hat x^{\mathrm f}_t)+\Sigma^{\mathrm f}$; since $\rho(A)<1$, $\Cov(x_t)\to\Sigma_x$ and $\E x_t=A^tm_0\to0$, so $\E d_td_t^\top\to D(\Sigma_x-\Sigma^{\mathrm f})D^\top$.
\end{proof}

\begin{proof}[Proof of \cref{cor:apx-unconditional-memory}]
For $M\succeq0$, $\det(I+M)\ge1+\tr M$, so with $M=(Q^y)^{-1/2}D\Sigma_xD^\top(Q^y)^{-1/2}$ the formula for $\mathfrak d_{\mathrm{cl}}$ in the proof of \cref{cor:apx-memory} gives $\tr((Q^y)^{-1}D\Sigma_xD^\top)\le\mathsf e_{\mathrm{cl}}$. Here $\bar Q=Q^y$, so $\mathsf N=\Delta$, $\mathsf l=1-\theta^0_{\mathrm b}$ and $P^-_{\mathrm b}\succeq Q^y$. Since $\Sigma^{\mathrm f}\preceq\Sigma_x$, $\Vert(P^-_{\mathrm b})^{-1/2}D\Sigma^{\mathrm f}D^\top(P^-_{\mathrm b})^{-1/2}\Vert\le\Vert M\Vert\le\mathsf e_{\mathrm{cl}}$. By the proof of \cref{cor:apx-memory}, $\rho(\bar A)<1$ and $\Sigma_y=\bar A\Sigma_y\bar A^\top+D\Sigma_xD^\top+Q^y$, so $\bar AC\Sigma^{\mathrm f}C^\top\bar A^\top\preceq\bar A\Sigma_y\bar A^\top\preceq\Sigma_y$ and $\Vert(P^-_{\mathrm b})^{-1/2}\bar AC(\Sigma^{\mathrm f})^{1/2}\Vert^2\le\mathsf k_y$. Hence the two cross terms of $\tilde\Delta$ have norm at most $(\mathsf k_y\mathsf e_{\mathrm{cl}})^{1/2}$ each and $\Vert\tilde{\mathsf N}\Vert\le\bar{\mathsf n}$. Moreover $\Sigma_y$ dominates the solution of $Y=\bar AY\bar A^\top+Q^y$, which dominates $P^-_{\mathrm b}$ by \cref{lem:apx-riccati}(d), so $1-\theta^0_{\mathrm b}=\lambda_{\min}(\tilde{Q^y})\ge\lambda_{\min}(Q^y)/\Vert\Sigma_y\Vert=1/\mathsf k_y$. By \cref{thm:apx-unconditional}(b), $\mathsf u\le1+\bar{\mathsf n}/(1-\theta_{\mathrm b})$ and $\mathsf w\le\bar{\mathsf n}/(1-\theta^0_{\mathrm b})$, so $\mathsf r\le\bar{\mathsf r}$. Without input $\E\Vert d_t\Vert^2_{(P^-_{\mathrm b})^{-1}}\le\E\Vert\delta_t\Vert^2_{(Q^y)^{-1}}$, whose upper limit is $\tr((Q^y)^{-1}D\Sigma_xD^\top)\le\mathsf e_{\mathrm{cl}}$ by the argument for the defect energy in the proof of \cref{cor:apx-memory}. Inserting these in \cref{eq:apx-unconditional}, with $\Vert\mathsf Y\Vert_{\mathrm F}\le\sqrt k\bar{\mathsf r}$ and the bound on the first term from \cref{thm:apx-unconditional}(c), gives \cref{eq:apx-unconditional-memory}. As $\mathfrak d_{\mathrm{cl}}\to0$, $\mathsf e_{\mathrm{cl}}=O(\mathfrak d_{\mathrm{cl}})$ and $\bar{\mathsf n}=O(\mathfrak d_{\mathrm{cl}}^{1/2})$ when $\mathsf k_y$ stays bounded, and then $\theta_{\mathrm b}\le\theta^0_{\mathrm b}\le1-1/\mathsf k_y$ bounds $\tau_{\mathrm b}$, $(1-\theta_{\mathrm b})^{-1}$ and $(1-\theta^0_{\mathrm b})^{-1}$, so $\bar{\mathsf r}^2=O(\mathfrak d_{\mathrm{cl}})$, and the right side is $O(\mathfrak d_{\mathrm{cl}})$ when also $k$ and $\mathsf j$ stay bounded.
\end{proof}

\begin{proof}[Proof of \cref{prop:apx-start}]
Write $P^-=P^-_{\mathrm b}$ and $\Vert v\Vert_{\mathrm b}=(\E\Vert v\Vert^2_{(P^-)^{-1}})^{1/2}$.

(a) \Cref{lem:apx-riccati}(c),(d),(e) with $\mathsf F=A$, $\mathsf J=HC$ and $N=Q\succ0$ give the solution $\Sigma^-\succeq Q$ of $\Sigma^-=\mathcal T_Q(\Sigma^-)$, its uniqueness among positive semidefinite solutions, and $\Sigma^-_t=\mathcal T_Q^t(\Sigma_0)\to\Sigma^-$. The map from $\Sigma^-_t$ to $\Sigma^{\mathrm f}_t$ in \cref{eq:apx-record-filter} is continuous, so $\Sigma^{\mathrm f}_t\to\Sigma^{\mathrm f}$. The conditional laws are those of \cref{lem:apx-separation}(b) in the form of \cref{rem:apx-separation-start}, and $P_{\mathrm e,t}=CA\Sigma^{\mathrm f}_{t-1}A^\top C^\top+Q^y\succeq Q^y\succ0$ for $t\ge1$. The equivalence stated after (c) is \cref{lem:apx-riccati}(e).

(b) The proofs of parts (a) and (b) of \cref{thm:apx-unconditional} use the fine process only through the identity \cref{eq:apx-fixed-point} and through $P_{\mathrm e}\succeq Q^y$. Both hold for the fixed point of (a): the derivation of \cref{eq:apx-fixed-point} uses $\Sigma^-=A\Sigma^{\mathrm f}A^\top+Q$ and the middle formula of \cref{eq:apx-record-filter} and nothing else.

(c) Both laws in $\Delta\mathcal F_t$ are Gaussian, so for $t\ge1$, with $e_t=C\hat x^-_t-\hat\eta^-_t$,
\begin{equation*}
\Delta\mathcal F_t=h\bigl((P^-)^{-1}P_{\mathrm e,t}\bigr)+\tfrac12\Vert e_t\Vert_{\mathrm b}^2,\qquad h\bigl((P^-)^{-1}P_{\mathrm e,t}\bigr)\to h\bigl((P^-)^{-1}P_{\mathrm e}\bigr)=\tfrac12\sum_{i=1}^kg(1+\mathsf y_i) .
\end{equation*}
The innovation $\iota_t=s_{t+1}-HC\hat x^-_t$ is, given $(s_{1:t},\omega)$, Gaussian with mean zero and the deterministic covariance $\Sigma_{\iota,t}=HP_{\mathrm e,t}H^\top+R$, so the innovations are independent with mean zero, and they are independent of $\omega$. With $L_{\mathrm e,t}=P_{\mathrm e,t}H^\top\Sigma_{\iota,t}^{-1}$ the derivation of \cref{eq:apx-error-recursion} gives
\begin{equation*}
e_{t+1}=\Psi'e_t+\bar A(L_{\mathrm e,t}-L)\iota_t+d_t,\qquad d_t=D\hat x^{\mathrm f}_t+Eu_t=\E[\delta_t\mid s_{1:t+1},\omega] .
\end{equation*}
By \cref{lem:apx-difference} with $P_1=P_{\mathrm e,t}$ and $P_2=P^-$, $\bar A(L_{\mathrm e,t}-L)=\Psi'\mathsf X_tH^\top\Sigma_{\iota,t}^{-1}$ with $\mathsf X_t=P_{\mathrm e,t}-P^-$, so the middle term has covariance $\Psi'\mathsf X_tJ_{\mathrm e,t}\mathsf X_t\Psi'^\top$ with $J_{\mathrm e,t}=H^\top\Sigma_{\iota,t}^{-1}H\preceq H^\top R^{-1}H$. Solving the recursion as in the proof of \cref{thm:apx-unconditional}(c) and using \cref{eq:apx-metric-contraction}, with $\mathsf Y_s=\tilde{\mathsf X}_s$,
\begin{equation*}
\Vert e_t\Vert_{\mathrm b}\le\theta_{\mathrm b}^{t/2}\Vert e_0\Vert_{\mathrm b}+\sum_{s<t}\theta_{\mathrm b}^{(t-1-s)/2}\Vert d_s\Vert_{\mathrm b}+\Bigl(\mathsf j\sum_{s<t}\theta_{\mathrm b}^{t-s}\Vert\mathsf Y_s\Vert_{\mathrm F}^2\Bigr)^{1/2} .
\end{equation*}
Inserted into the first display, this is \cref{eq:apx-start-round}. The first term tends to zero because $\Vert e_0\Vert_{\mathrm b}<\infty$. Every $\Vert d_s\Vert_{\mathrm b}$ is finite, so the upper limit of the second term is at most $\tau_{\mathrm b}\bar d_{\mathrm b}$. Since $\mathsf Y_s\to\mathsf Y$ by (a) and $\theta_{\mathrm b}<1$, the sum in the third term converges to $\theta_{\mathrm b}\Vert\mathsf Y\Vert_{\mathrm F}^2/(1-\theta_{\mathrm b})$. This gives \cref{eq:apx-unconditional}. The inequality $\Vert d_t\Vert_{\mathrm b}^2\le\E\Vert\delta_t\Vert^2_{\bar Q^{-1}}$ holds as before, because $u_t$ is measurable in $(s_{1:t},\omega)$, a conditional expectation does not increase $\Vert\cdot\Vert_{\mathrm b}$, and $(P^-)^{-1}\preceq\bar Q^{-1}$.
\end{proof}

\begin{proposition}[Finite horizon without stationarity]\label{prop:apx-posterior-finite}
Let the meta-agent run the time-varying filter of the declared model from the declared initial law: $\hat P^-_0=\hat\Sigma_0$, $\hat\eta^-_0=\hat m_0$, $L_t=\hat P^-_tH^\top(H\hat P^-_tH^\top+R)^{-1}$, $\hat P^-_{t+1}=\bar A\varphi_R(\hat P^-_t)\bar A^\top+\bar Q$ and $\hat\eta^-_{t+1}=\bar A(\hat\eta^-_t+L_t(s_{t+1}-H\hat\eta^-_t))+\bar Bu_t$. Let the policy be $u_t=\pi_t(s_{1:t},\omega)$ with $\E\lvert u_t\rvert^2<\infty$. Without any hypothesis of stability, stationarity or detectability, for every $t\ge0$,
\begin{equation}\label{eq:apx-posterior-finite}
\E\KL\Bigl(P\bigl(y_t\in\cdot\mid s_{1:t},\omega\bigr)\Big\Vert\Ncal\bigl(\hat\eta^-_t,\hat P^-_t\bigr)\Bigr)\le\KL\bigl(P^{\pi'}\Vert\hat P^{\pi'}\bigr),
\end{equation}
where the right side is the path relative entropy of \cref{thm:apx-path} at horizon $T=t$ for the causal coarse policy $\pi'_j(y_{0:j},\omega')=\pi_j(Hy_0+v_0,\dots,Hy_{j-1}+v_{j-1},\omega)$ with randomization $\omega'=(\omega,v_{0:t-1})$. At $t=0$ the two sides are equal.
\end{proposition}

\begin{proof}
Let $\mathbb P$ and $\hat{\mathbb P}$ be the laws of $(\omega,s_{1:t},y_t)$ under the fine process and under the declared model read through the same channel and driven by the same maps. By \cref{lem:apx-separation}(b) applied to the declared model, with $(A,B,Q,C)$ replaced by $(\bar A,\bar B,\bar Q,I)$, the conditional law of $y_t$ given $(\omega,s_{1:t})$ under $\hat{\mathbb P}$ is $\Ncal(\hat\eta^-_t,\hat P^-_t)$. The chain rule gives
\begin{equation*}
\KL\bigl(\mathbb P\Vert\hat{\mathbb P}\bigr)=\KL\bigl(\mathbb P_{(\omega,s_{1:t})}\Vert\hat{\mathbb P}_{(\omega,s_{1:t})}\bigr)+\E\KL\Bigl(P\bigl(y_t\in\cdot\mid s_{1:t},\omega\bigr)\Big\Vert\Ncal\bigl(\hat\eta^-_t,\hat P^-_t\bigr)\Bigr),
\end{equation*}
and the first term on the right is nonnegative. The vector $(\omega,s_{1:t},y_t)$ is the same measurable function of $(\omega',y_{0:t})$ under both laws, $\omega'$ has the same law under both and is independent of $(x_0,\mathfrak n)$ and of $(\hat y_0,\hat w)$, and $\pi'$ is a causal coarse policy. Relative entropy does not increase under a measurable map, so $\KL(\mathbb P\Vert\hat{\mathbb P})\le\KL(P^{\pi'}\Vert\hat P^{\pi'})$. At $t=0$ both sides of \cref{eq:apx-posterior-finite} are $\KL(P_{y_0}\Vert\hat P_0)$.
\end{proof}

\begin{proof}[Proof of \cref{prop:apx-score}]
Here $\Lambda^\star\succeq0$ is symmetric, the offsets $c_t$ and the targets $o^\star_t$ are deterministic, $G_t=\frac12\E^\pi\Vert y_t-o^\star_t\Vert^2_{\Lambda^\star}$ and $\hat G_t$ is the same expression under the declared model. Under the fine process $y_{t+1}=\bar A_Ky_t+\bar Bc_t+\delta_t+C\mathfrak n_t$. Let $\zeta_t=(Q^y)^{-1/2}C\mathfrak n_t$, which are independent standard normal vectors, and define on the same probability space $\hat y_{t+1}=\bar A_K\hat y_t+\bar Bc_t+\bar Q^{1/2}\zeta_t$, where $\hat y_0$ is drawn from the conditional law of an optimal coupling of $(P_{y_0},\hat P_0)$ given $y_0$, independently of everything else. Then $\hat y_0$ is independent of $(\zeta_t)$, so $\hat y$ has the law of the declared model under the same feedback, and $\Vert y_0-\hat y_0\Vert_{L^2}=w_0$. The difference $\Delta_t=y_t-\hat y_t$ obeys $\Delta_{t+1}=\bar A_K\Delta_t+\delta_t+\Gamma\zeta_t$ with $\Gamma=(Q^y)^{1/2}-\bar Q^{1/2}$, hence
\begin{equation*}
\Delta_t=\bar A_K^t\Delta_0+\sum_{s<t}\bar A_K^{t-1-s}\delta_s+\sum_{s<t}\bar A_K^{t-1-s}\Gamma\zeta_s .
\end{equation*}
The last sum has squared $L^2$ norm $\sum_{s<t}\Vert\bar A_K^{t-1-s}\Gamma\Vert_{\mathrm F}^2\le\beta^2\sum_{j<t}\Vert\bar A_K^j\Vert^2$, and the triangle inequality in $L^2$ gives $W_2(P^\pi_{y_t},\hat P^\pi_{y_t})\le\Vert\Delta_t\Vert_{L^2}\le w_t$.

For the score, $G_t^{1/2}$ and $\hat G_t^{1/2}$ are the $L^2$ norms of $(\Lambda^\star/2)^{1/2}(y_t-o^\star_t)$ and $(\Lambda^\star/2)^{1/2}(\hat y_t-o^\star_t)$, so $\lvert G_t^{1/2}-\hat G_t^{1/2}\rvert\le(\Vert\Lambda^\star\Vert/2)^{1/2}\Vert\Delta_t\Vert_{L^2}\le a_t$ and $\lvert G_t-\hat G_t\rvert=\lvert G_t^{1/2}-\hat G_t^{1/2}\rvert(G_t^{1/2}+\hat G_t^{1/2})\le a_t(2\hat G_t^{1/2}+a_t)$. Summing over $t$ bounds $\lvert G(\pi)-\hat G(\pi)\rvert$.

For state feedback $u_t=k_t(y_t)$ with $\lvert k_t(y)-k_t(y')\rvert\le\kappa\lvert y-y'\rvert$, the same coupling gives $\Delta_{t+1}=\bar A\Delta_t+\bar B(k_t(y_t)-k_t(\hat y_t))+\delta_t+\Gamma\zeta_t$, so with $\ell=\Vert\bar A\Vert+\kappa\Vert\bar B\Vert$ and $\Vert\Gamma\zeta_t\Vert_{L^2}=\beta$, $\Vert\Delta_{t+1}\Vert_{L^2}\le\ell\Vert\Delta_t\Vert_{L^2}+\Vert\delta_t\Vert_{L^2}+\beta$, and induction gives $\Vert\Delta_t\Vert_{L^2}\le\ell^tw_0+\sum_{s<t}\ell^{t-1-s}(\Vert\delta_s\Vert_{L^2}+\beta)$.
\end{proof}

\begin{proof}[Proof of \cref{prop:apx-softmax}]
Here $\gamma>0$ and $E(\pi)>0$ for every $\pi\in\Pi$. Let $f=\hat G-G$ and $R=\gamma\operatorname{osc}$. Then $\hat p/p=e^{-\gamma f}/\E_pe^{-\gamma f}$.

Log-ratio. Since $-\gamma\max f\le\ln\E_pe^{-\gamma f}\le-\gamma\min f$, $\ln(\hat p(\pi)/p(\pi))=-\gamma f(\pi)-\ln\E_pe^{-\gamma f}$ lies between $-\gamma(f(\pi)-\min f)$ and $\gamma(\max f-f(\pi))$, hence in $[-R,R]$.

Relative entropy. $\KL(p\Vert\hat p)=\gamma\E_pf+\ln\E_pe^{-\gamma f}=\psi(-\gamma)$ with $\psi(s)=\ln\E_pe^{s(f-\E_pf)}$. One has $\psi(0)=0$, $\psi'(0)=0$, and $\psi''(s)$ is the variance of $f$ under the law proportional to $pe^{sf}$, which is at most $\operatorname{osc}^2/4$ because $f$ takes values in an interval of length $\operatorname{osc}$. Taylor's formula gives $\psi(-\gamma)\le\gamma^2\operatorname{osc}^2/8$.

Total variation. The distance is $\max_{S\subseteq\Pi}(\hat p(S)-p(S))$. With the weights $w=e^{-\gamma(f-\max f)}\in[1,e^R]$ and $a=p(S)$,
\begin{equation*}
\hat p(S)=\frac{\sum_Spw}{\sum_Spw+\sum_{\Pi\setminus S}pw}\le\frac{ae^R}{ae^R+1-a},
\end{equation*}
because the ratio increases with the first sum and decreases with the second. The function $a\mapsto ae^R/(ae^R+1-a)-a$ on $[0,1]$ is maximal at $a=1/(1+e^{R/2})$, where it equals $(e^{R/2}-1)/(e^{R/2}+1)=\tanh(R/4)$. The value is attained: for two policies with equal habits, $\gamma(G(\pi_1)-G(\pi_2))=R/2$ and $\gamma(\hat G(\pi_1)-\hat G(\pi_2))=-R/2$, one has $p=(a,1-a)$ and $\hat p=(1-a,a)$ with $a=1/(1+e^{R/2})$, and $1-2a=\tanh(R/4)$.
\end{proof}

\begin{proof}[Proof of \cref{thm:apx-levels}]
(a) With $C^{(0)}=I$, $D^{\mathrm{tot}}_0=0$ and $E^{\mathrm{tot}}_0=0$,
\begin{equation*}
D^{\mathrm{tot}}_{l+1}=C_{l+1}\bigl(D^{\mathrm{tot}}_l+\bar A_lC^{(l)}\bigr)-\bar A_{l+1}C_{l+1}C^{(l)}=C_{l+1}D^{\mathrm{tot}}_l+D_{l+1}C^{(l)},
\end{equation*}
and $E^{\mathrm{tot}}_{l+1}=C_{l+1}(E^{\mathrm{tot}}_l+\bar B_l)-\bar B_{l+1}=C_{l+1}E^{\mathrm{tot}}_l+E_{l+1}$.

(b) By (a) and $y^{(l)}_t=C^{(l)}x_t$, $D^{\mathrm{tot}}_{l+1}x_t+E^{\mathrm{tot}}_{l+1}u_t=C_{l+1}(D^{\mathrm{tot}}_lx_t+E^{\mathrm{tot}}_lu_t)+D_{l+1}y^{(l)}_t+E_{l+1}u_t$, and the triangle inequality in $L^2$ gives the recursion. A policy that is causal in the top-level variable is a causal coarse policy for $C=C^{(l)}$ at every level $l$, since the top-level variable is a linear image of $y^{(l)}$. If $C^{(l)}$ has full row rank and $C^{(l)}\Sigma_0C^{(l)\top}\succ0$, \cref{thm:apx-path} applies to the retained variable $y^{(l)}$ and the declared model $(\bar A_l,\bar B_l,\bar Q_l)$, with $\delta_t=D^{\mathrm{tot}}_lx_t+E^{\mathrm{tot}}_lu_t$; dropping $\mathfrak d_{\mathrm{dyn},t}\ge0$ and using $\E\lvert\delta_t\rvert^2\le\mathfrak e_l^2$ gives \cref{eq:apx-uniform} with $Q^y=C^{(l)}QC^{(l)\top}$, $\bar Q=\bar Q_l$ and $\bar{\mathfrak d}=\mathfrak e_l$.

(c) Let $\mu_l$ and $\nu_l$ be the laws of $y^{(l)}_t$ under the fine process and under the level-$l$ model, both of finite second moment, and let $C_{l+1}\mu_l$ denote the image law. Then $\mu_{l+1}=C_{l+1}\mu_l$, the image of a coupling of $(\mu_l,\nu_l)$ under $C_{l+1}$ applied to both entries is a coupling of the images, so $W_2(C_{l+1}\mu_l,C_{l+1}\nu_l)\le\Vert C_{l+1}\Vert\varepsilon_l$, and the triangle inequality for $W_2$ gives $\varepsilon_{l+1}\le\Vert C_{l+1}\Vert\varepsilon_l+W_2(C_{l+1}\nu_l,\nu_{l+1})$. The statement is metric and uses nothing about the models beyond finite second moments.
\end{proof}

\begin{proof}[Proof of \cref{thm:apx-h5a}]
Reduction to the setting of \cref{prop:apx-start}. In the canonical partition no exterior variable reads an internal state, and (a2) gives $A_{\sE\sS}=0$ and $A_{\sE\sA}=\mathsf L_{\mathrm E}\mathsf P_{\mathrm A}$, so $X_{\sE,t+1}=A_{\sE\sE}X_{\sE,t}+\mathsf L_{\mathrm E}a_{C,t}+\Xi^{\sE}_t$ with $Q^{\Xi}_{\sE\sE}\succ0$; this is the fine process of \cref{def:apx-setting} with the identifications stated before the theorem, and $\mathsf P_{\mathrm E}$ has full row rank. Since $\Lambda^s=\mathcal H^\top R_{\sS}^{-1}\mathcal H\succ0$, the matrix $\mathcal H$ of (S) has full column rank, $(\bar F,\Lambda^s)$ is detectable because $\Lambda^s$ is invertible, and the steady-state covariances of the declared model exist and are unique (\cref{lem:apx-riccati}(c),(d)). For $H=R=\Lambda^s$, $\varphi_R(P)=(P^{-1}+\Lambda^s)^{-1}$ for $P\succ0$ and $PH^\top(HPH^\top+R)^{-1}=P(\Lambda^sP+I)^{-1}=(P^{-1}+\Lambda^s)^{-1}$, so $P^-_{\mathrm b}=\bar P^-$, $P_{\mathrm b}=L=\bar P$, $\Psi'=\bar F(I-\bar P\Lambda^s)=\mathsf K$ and $H^\top R^{-1}H=\Lambda^s$.

By (S), $X_{\sS,t+1}=\mathcal H\eta_{C,t}+\Xi^{\sS}_t$ and $s_{C,t+1}=\Lambda^s\eta_{C,t}+\bar v_t$ with $\bar v_t=\mathcal H^\top R_{\sS}^{-1}\Xi^{\sS}_t\sim\Ncal(0,\Lambda^s)$, and
\begin{equation*}
\begin{gathered}
X_{\sS,t+1}=\mathcal H(\Lambda^s)^{-1}s_{C,t+1}+\mathfrak r^\perp_t,\qquad\mathfrak r^\perp_t=\bigl(I-\mathcal H(\Lambda^s)^{-1}\mathcal H^\top R_{\sS}^{-1}\bigr)\Xi^{\sS}_t,\\
\Cov(\mathfrak r^\perp_t,\bar v_t)=\bigl(I-\mathcal H(\Lambda^s)^{-1}\mathcal H^\top R_{\sS}^{-1}\bigr)\mathcal H=0,
\end{gathered}
\end{equation*}
so the jointly Gaussian vectors $\mathfrak r^\perp_t$ and $\bar v_t$ are independent. The internal and active states read collective variables only, and the sensory channels read no interior or active variable. Let $\omega=(X_{\sI,0},X_{\sS,0},X_{\sA,0},(\Xi^{\sI}_j,\Xi^{\sA}_j,\mathfrak r^\perp_j)_{j\ge0})$. By induction on $t$, the sensory states at rounds $0,\dots,t$ and the internal and active states at rounds $0,\dots,t+1$ are measurable functions of $(s_{C,1:t},\omega)$. With per-channel noise and an initial state that is independent of the noise, $(\Xi^{\sE},\bar v)$ is independent of $(X_{\sE,0},\omega)$, and the law of $X_{\sE,0}$ given $\omega$ is its law given the initial collective variables, $\Ncal(m_0,\Sigma_0)$ by (H-prot). Hence $a_{C,t}=\pi_t(s_{C,1:t},\omega)$ for measurable maps $\pi_t$, and $\E\lvert a_{C,t}\rvert^2<\infty$ because the joint process has finite second moments from a start that has them. This is a record-measurable policy in the sense of \cref{sec:apx} with a start of the kind of \cref{rem:apx-separation-start}, and \cref{lem:apx-separation}(b) applies in that form.

The fixed point $\Sigma^-$. Let $\mathcal T_{\sE}(P)=A_{\sE\sE}\varphi(P)A_{\sE\sE}^\top+Q^{\Xi}_{\sE\sE}$ in the notation of \cref{lem:apx-riccati} with $\mathsf F=A_{\sE\sE}$, $\mathsf J=\Lambda^s\mathsf P_{\mathrm E}$, $R=\Lambda^s$ and $N=Q^{\Xi}_{\sE\sE}$, so that \eqref{eq:apx-record-filter} reads $\Sigma^-_{t+1}=\mathcal T_{\sE}(\Sigma^-_t)$. Consider the joint process from the start $X_0=0$, which is a start of (H-prot) with $m_0=0$ and $\Sigma_0=0$. By \cref{lem:apx-separation}(b) the covariance of $X_{\sE,t}$ given $(s_{C,1:t},\omega)$ is the deterministic matrix $\mathcal T_{\sE}^t(0)$, and a deterministic conditional covariance is at most the covariance, so $\mathcal T_{\sE}^t(0)\preceq\Cov(X_{\sE,t})$; and $\Cov(X_t)=\sum_{j<t}A^jQ^{\Xi}A^{j\top}$ is bounded in $t$ because $\rho(A)<1$. By \cref{lem:apx-riccati}(e) the fixed point $\Sigma^-$ exists, it is unique, $(A_{\sE\sE},\Lambda^s\mathsf P_{\mathrm E})$ is detectable, and $\Sigma^-_t\to\Sigma^-$ from every $\Sigma_0\succeq0$; no property of $A_{\sE\sE}$ alone was used. Also $\mathsf P_{\mathrm E}\Sigma^-\mathsf P_{\mathrm E}^\top\succeq Q^y\succ0$, and \cref{prop:apx-start} applies.

(i) The proof of \cref{thm:met-h5a}(i) uses (H-suff) and $\rho(A)<1$ and nothing of (H-phys). If $X_{\sE,0}$ is independent of the initial collective variables, of which $(X_{\sI},\mathfrak r_{\mathrm S},\mathfrak r_{\mathrm A},b_C)_0$ is a function, then (a3) holds.

(ii) The proof of \cref{thm:met-h5a}(ii) uses only the form of the update \eqref{eq:met-member-filter}, a matrix $\mathcal W$ that preserves sums and coefficients common to the members, and it is stated there for any fixed matrices in place of $(\bar F,\bar P,\bar B,\Lambda^s)$.

(iii) Write $\mathfrak z_0$ for the initial collective variables and $\omega'$ for the noises of the collective variables, so that $(\mathfrak z_0,\omega')$ is a function of $(\omega,\bar v)$ and conversely. Since $s_{C,j+1}=\Lambda^s\eta_{C,j}+\bar v_j$, $a_{C,t}=\pi_t(\Lambda^s\eta_{C,0}+\bar v_0,\dots,\Lambda^s\eta_{C,t-1}+\bar v_{t-1},\omega)=\pi'_t(\eta_{C,0:t},\omega')$ for maps $\pi'_t$ that depend on $\mathfrak z_0$. Given $\mathfrak z_0$, the exterior starts from $\Ncal(m_0,\Sigma_0)$ with $\mathsf P_{\mathrm E}\Sigma_0\mathsf P_{\mathrm E}^\top\succ0$, the randomization $\omega'$ is independent of $(X_{\sE,0},\Xi^{\sE})$ and has the law it has under the declared model, and $\E[\lvert a_{C,t}\rvert^2\mid\mathfrak z_0]<\infty$ almost surely. So, given $\mathfrak z_0$, $\pi'$ is a causal coarse policy, as in the proof of \cref{prop:apx-posterior-finite}, and \cref{thm:apx-path} gives \eqref{eq:apx-path} for the conditional laws, with first term $\KL(\Ncal(\mathsf P_{\mathrm E}m_0,\mathsf P_{\mathrm E}\Sigma_0\mathsf P_{\mathrm E}^\top)\Vert\hat P_0)$, with the conditional expectation of the defect energy, and with covariances $P_t$ that do not depend on $\mathfrak z_0$. The two laws give $\mathfrak z_0$ the same marginal, so by the chain rule their relative entropy is the expectation over $\mathfrak z_0$ of the relative entropy of the conditional laws, which is the formula of the statement; its first term is finite because $\E\lvert m_0\rvert^2<\infty$. For the stationary start, $X_{\sE,0}$ has the law of $A_{\sE\sE}X'_{\sE}+\mathsf L_{\mathrm E}a'_C+\Xi'$ for a state $X'$ one round earlier and a noise $\Xi'\sim\Ncal(0,Q^{\Xi}_{\sE\sE})$ that is independent of $X'$ and, with per-channel noise, of the collective variables $\mathfrak z_0$ of the current round, which are functions of $X'$ and of the noises of the collective channels; hence the covariance of $X_{\sE,0}$ given $\mathfrak z_0$ is at least its covariance given $(X',\mathfrak z_0)$, which is $Q^{\Xi}_{\sE\sE}$. The defect $\delta_t=DX_{\sE,t}+E\mathsf P_{\mathrm A}X_{\sA,t}$ is a linear function of $X_t$; since $\rho(A)<1$ and the start has finite second moments, $\E X_t\to0$ and $\Cov(X_t)$ converges to the stationary covariance, so $\E\delta_t\delta_t^\top\to\Sigma_\delta$ and $\E\Vert\delta_t\Vert^2_{\bar Q^{-1}}\to\tr(\bar Q^{-1}\Sigma_\delta)$. Dropping $\mathfrak d_{\mathrm{dyn},t}\ge0$ in \eqref{eq:apx-path}, dividing by $T$ and taking the limit of the averages gives \eqref{eq:apx-h5a-path}. If $D=0$ and $E=0$ then $\delta_t=0$, and $\mathsf P_{\mathrm E}A_{\sE\sE}P_t=\bar F\mathsf P_{\mathrm E}P_t=0$ because $\mathsf P_{\mathrm E}P_t=0$ (\cref{lem:apx-separation}(a) with $C=\mathsf P_{\mathrm E}$), so $\mathfrak d_{\mathrm{dyn},t}=0$; with $\bar Q=Q^y$ also $h(\bar Q^{-1}Q^y)=0$, and every summand of \eqref{eq:apx-path} vanishes.

(iv) By \cref{lem:apx-separation}(b) in the form of \cref{rem:apx-separation-start}, the law of $X_{\sE,t}$ given $(s_{C,1:t},\omega)$ is $\Ncal(\hat x^-_t,\Sigma^-_t)$, where $\hat x^-_0=m_0$, a function of the initial collective variables, and $\hat x^-_{j+1}=A_{\sE\sE}\hat x^{\mathrm f}_j+\mathsf L_{\mathrm E}a_{C,j}$, with $\hat x^{\mathrm f}_j$ an affine function of $\hat x^-_j$ and $s_{C,j+1}$ with deterministic coefficients, so $\hat x^-_t$ is a function of $m_0$ and $(s_{C,1:t},a_{C,0:t-1})$. Let $\mathcal B_t$ be generated by $(s_{C,1:t},a_{C,0:t-1})$ and $\mathcal G^C_t$ by the collective variables at rounds $0,\dots,t$. Then $\mathcal B_t\subseteq\mathcal G^C_t\subseteq\sigma(s_{C,1:t},\omega)$ by the induction above, $\hat x^-_t$ is measurable in $\mathcal G^C_t$, and a conditional law given the larger of two nested $\sigma$-fields that is measurable in the smaller is the conditional law given the smaller. This is the first claim; when $m_0$ is deterministic, $\hat x^-_t$ is measurable in $\mathcal B_t$ and the same argument gives the law given the record. The law of $\eta_{C,t}=\mathsf P_{\mathrm E}X_{\sE,t}$ follows, and the convergence of the covariances was shown above.

By (ii) with the computation noises written out, $y^\eta_{C,t+1}=\mathsf Ky^\eta_{C,t}+\bar F\bar Ps_{C,t}+\bar By^\alpha_{C,t}+\bar\nu^\mu_t$ and $y^\alpha_{C,t}=a_{C,t-1}+\bar\nu^\alpha_{t-1}$ for $t\ge1$, and the steady-state filter of the declared model is $\hat\eta^-_t=\mathsf K\hat\eta^-_{t-1}+\bar F\bar Ps_{C,t}+\bar Ba_{C,t-1}$, which is the recursion of \cref{sec:apx} because $\bar A(\hat\eta^-_{t-1}+L(s_t-H\hat\eta^-_{t-1}))=\mathsf K\hat\eta^-_{t-1}+\bar F\bar Ps_{C,t}$. Subtracting, $\varepsilon^{\mathrm c}_t=y^\eta_{C,t+1}-\hat\eta^-_t$ obeys $\varepsilon^{\mathrm c}_t=\mathsf K\varepsilon^{\mathrm c}_{t-1}+\nu^{\mathrm c}_t$ for $t\ge1$, and $\varepsilon^{\mathrm c}_0=0$ for the start $\hat\eta^-_0=y^\eta_{C,1}$, a function of $\omega$ with finite second moment. Both laws in $\Delta\mathcal F^C_t$ are Gaussian, and $P_{\mathrm e,t}\succeq Q^y\succ0$ for $t\ge1$, which gives \eqref{eq:apx-h5a-split} as in \cref{eq:apx-excess-split}, and $\mathsf P_{\mathrm E}\hat x^-_t-y^\eta_{C,t+1}=e_t-\varepsilon^{\mathrm c}_t$ with $e_t=\mathsf P_{\mathrm E}\hat x^-_t-\hat\eta^-_t$. Write $\Vert v\Vert_{\mathrm b}=(\E\Vert v\Vert^2_{(\bar P^-)^{-1}})^{1/2}$. The proof of \cref{prop:apx-start}(c) bounds $e_t$ through the triangle inequality for $\Vert\cdot\Vert_{\mathrm b}$ applied to the solution of the error recursion with the gains $L_{\mathrm e,t}$; it uses the initial value of the declared filter only through $\Vert e_0\Vert_{\mathrm b}<\infty$, which holds here because $e_0=\mathsf P_{\mathrm E}m_0-y^\eta_{C,1}$, and gives $\limsup_t\Vert e_t\Vert_{\mathrm b}\le\tau_{\mathrm b}\bar d_{\mathrm b}+(\theta_{\mathrm b}\mathsf j/(1-\theta_{\mathrm b}))^{1/2}\Vert\mathsf Y\Vert_{\mathrm F}$ for every $\bar d_{\mathrm b}\ge\limsup_t\Vert d_t\Vert_{\mathrm b}$. Here $d_t=\E[\delta_t\mid s_{C,1:t+1},\omega]=D\hat x^{\mathrm f}_t+Ea_{C,t}$, because $a_{C,t}$ is measurable in $(s_{C,1:t},\omega)$; it is a function of the collective variables through round $t+1$, hence also the mean of $\delta_t$ given those variables. Since $\delta_t-d_t=D(X_{\sE,t}-\hat x^{\mathrm f}_t)$ has conditional covariance $D\Sigma^{\mathrm f}_tD^\top$ and conditional mean zero, $\E d_td_t^\top=\E\delta_t\delta_t^\top-D\Sigma^{\mathrm f}_tD^\top\to\Sigma_\delta-D\Sigma^{\mathrm f}D^\top$, which gives the stated $\bar d_{\mathrm b}$ as a limit; it is at most $\tr((\bar P^-)^{-1}\Sigma_\delta)\le\tr(\bar Q^{-1}\Sigma_\delta)$ because $\bar P^-\succeq\bar Q$. For the response, $\varepsilon^{\mathrm c}_t=\sum_{l<t}\mathsf K^l\nu^{\mathrm c}_{t-l}$, and by \cref{eq:apx-metric-contraction} in the form $\Vert(\bar P^-)^{-1/2}\mathsf K^lv\Vert\le\theta_{\mathrm b}^{l/2}\Vert(\bar P^-)^{-1/2}v\Vert$, $\Vert\varepsilon^{\mathrm c}_t\Vert_{\mathrm b}\le\sum_{l<t}\theta_{\mathrm b}^{l/2}\Vert\nu^{\mathrm c}_{t-l}\Vert_{\mathrm b}\le\tau_{\mathrm b}\bar\nu_{\mathrm b}$: the noises $\bar\nu^\mu_t$ and $\bar\nu^\alpha_{t-1}$ belong to different rounds and are independent, so $\nu^{\mathrm c}_t$ has covariance $\Sigma_{\bar\nu^\mu}+\bar B\Sigma_{\bar\nu^\alpha}\bar B^\top$ and $\Vert\nu^{\mathrm c}_t\Vert_{\mathrm b}=\bar\nu_{\mathrm b}$ for $t\ge1$. The triangle inequality for $\Vert\cdot\Vert_{\mathrm b}$ gives the first line of \eqref{eq:apx-h5a}, and with \eqref{eq:apx-h5a-split}, the convergence $P_{\mathrm e,t}\to P_{\mathrm e}$ and the eigenvalues $1+\mathsf y_i$ of $(\bar P^-)^{-1}P_{\mathrm e}$ the second. Parts (a) and (b) of \cref{thm:apx-unconditional} concern the two filter equations only (\cref{prop:apx-start}(b)). If $D=0$, $E=0$ and $\bar Q=Q^y$, then $\Delta=0$ and $\mathsf N=0$, so $\mathsf w=0$ and $\mathsf Y=0$ by \cref{eq:apx-enclosure}, and $\Sigma_\delta=0$ and $d_t=0$. In that case $\mathsf P_{\mathrm E}A_{\sE\sE}\Sigma^{\mathrm f}_tA_{\sE\sE}^\top\mathsf P_{\mathrm E}^\top=\bar F\mathsf P_{\mathrm E}\Sigma^{\mathrm f}_t\mathsf P_{\mathrm E}^\top\bar F^\top=\bar F\varphi_R(P_{\mathrm e,t})\bar F^\top$, so $P_{\mathrm e,t+1}=\bar F\varphi_R(P_{\mathrm e,t})\bar F^\top+\bar Q$, the covariance recursion of the declared filter, of which $\bar P^-$ is the fixed point; if $P_{\mathrm e,0}=\bar P^-$ then $P_{\mathrm e,t}=\bar P^-$ and $L_{\mathrm e,t}=L$ for all $t$, and the error recursion reads $e_{t+1}=\mathsf Ke_t$ with $e_0=\mathsf P_{\mathrm E}m_0-y^\eta_{C,1}$. The deviation recursion in the proof of \cref{thm:met-h5a}(iv) is obtained by subtracting the average recursion of (ii) from the member recursions and involves no property of the exterior.

(v) Interventions shift means and leave covariances unchanged, so with a common $P$ each member's expected free energy over acts is \eqref{eq:agy-act-efe} with $\bar F$ for $F$ and the effects $\bar{\mathsf M}_j$, as in the proof of \cref{thm:met-h5a}(v). The proof of \cref{thm:agy-joint-acts}(b) expands $\frac12\Vert\mathsf z_i+\sum_j\mathsf M_ju_j\Vert^2_{\Lambda^\star}$ and uses the effect models only through their being common to the members and to the group model, so it applies to the declared effects, and the rest of the proof of \cref{thm:met-h5a}(v) gives the group posterior $\hat{\mathbb P}_C$ with the tilts $\mathcal T_i$, whose potential is $\hat G_C+\sum_i\mathcal T_i$ up to a constant. Under the law of the joint process, given the paths of the collective variables the current exterior has mean $\hat x^-$ and a covariance that does not depend on $u$, and $\eta_C'=\mathsf P_{\mathrm E}(A_{\sE\sE}X_{\sE}+\mathsf L_{\mathrm E}a_C(u)+\Xi^{\sE})$, so $o^\star$ is Gaussian with mean $\mathsf H_\star\mathsf P_{\mathrm E}(A_{\sE\sE}\hat x^-+\mathsf L_{\mathrm E}a_C(u))$ and a covariance independent of $u$, and minus its expected log preference is $G^\star_C(u)$ plus a constant, by the computation in the first paragraph of the proof of \cref{thm:agy-joint-acts}. Let $\hat\mu(u)$ and $\mu^\star(u)$ be the vectors whose squared norms are $2\hat G_C(u)$ and $2G^\star_C(u)$. Since $\mathsf P_{\mathrm E}A_{\sE\sE}=D+\bar F\mathsf P_{\mathrm E}$ and $\mathsf P_{\mathrm E}\mathsf L_{\mathrm E}=E+\bar B$,
\begin{equation*}
\mu^\star(u)-\hat\mu(u)=\mathsf H_\star\bigl[D\hat x^-+Ea_C(u)+\bar F\bigl(\mathsf P_{\mathrm E}\hat x^--m_C\bigr)\bigr],\qquad\Vert\mu^\star(u)-\hat\mu(u)\Vert_{\Lambda^\star}=\sqrt2\mathsf a(u),
\end{equation*}
and $\lvert G^\star_C-\hat G_C\rvert=\frac12\lvert\Vert\mu^\star\Vert_{\Lambda^\star}-\Vert\hat\mu\Vert_{\Lambda^\star}\rvert(\Vert\mu^\star\Vert_{\Lambda^\star}+\Vert\hat\mu\Vert_{\Lambda^\star})\le\frac12\sqrt2\mathsf a(2\Vert\hat\mu\Vert_{\Lambda^\star}+\sqrt2\mathsf a)=\mathsf a(2\hat G_C^{1/2}+\mathsf a)$. \Cref{prop:apx-softmax} with the tilted product habits as the common habits of the two laws and $\operatorname{osc}\le2\max_u\lvert\hat G_C-G^\star_C\rvert$ gives the bound on the total variation distance. The last sentence of (v) is the first line of \eqref{eq:apx-h5a}.
\end{proof}
